\documentclass[11pt,reqno]{amsart}
\usepackage{amsmath,amssymb,amsthm}
\usepackage[margin=1in]{geometry}
\usepackage{booktabs}
\usepackage{array}
\usepackage[colorlinks=true,linkcolor=blue,citecolor=blue,urlcolor=blue]{hyperref}

\theoremstyle{plain}
\newtheorem{theorem}{Theorem}[section]
\newtheorem{proposition}[theorem]{Proposition}
\newtheorem{lemma}[theorem]{Lemma}
\newtheorem{corollary}[theorem]{Corollary}
\newtheorem{conjecture}[theorem]{Conjecture}
\theoremstyle{definition}
\newtheorem{definition}[theorem]{Definition}
\newtheorem{spec}[theorem]{Specification}
\theoremstyle{remark}
\newtheorem{remark}[theorem]{Remark}
\newtheorem{reading}[theorem]{Reading}

\newcommand{\R}{\mathbb{R}}
\newcommand{\C}{\mathbb{C}}
\newcommand{\Z}{\mathbb{Z}}
\newcommand{\Q}{\mathbb{Q}}
\newcommand{\F}{\mathbb{F}}
\newcommand{\Sph}{\mathbb{S}}
\newcommand{\THH}{\mathrm{THH}}
\newcommand{\THR}{\mathrm{THR}}
\newcommand{\TP}{\mathrm{TP}}
\newcommand{\HZ}{H\Z}
\newcommand{\ku}{\mathrm{ku}}
\newcommand{\KU}{\mathrm{KU}}
\newcommand{\Xbar}{\overline{X}}
\newcommand{\prodreg}{\sideset{}{^{\mathrm{reg}}}\prod}
\newcommand{\grade}[1]{\textup{[\textsc{#1}]}}
\DeclareMathOperator{\Lam}{\Lambda}
\DeclareMathOperator{\Hom}{Hom}
\DeclareMathOperator{\Ext}{Ext}
\DeclareMathOperator{\Aut}{Aut}
\DeclareMathOperator{\fib}{fib}
\DeclareMathOperator{\colim}{colim}
\DeclareMathOperator{\res}{res}
\DeclareMathOperator{\MAX}{MAX}
\DeclareMathOperator{\MIN}{MIN}
\DeclareMathOperator{\erfc}{erfc}
\emergencystretch=1.5em
\raggedbottom

% DRAFT assembled 2026-06-11 from the research/proof/ document set (17--37) and the
% research/work/ worker records cited at each statement. Every labeled claim carries a
% grading tag; numerical claims are reproducible via the cited witness scripts in
% scripts/numerics/witness/ (all exit-code gated, fail-closed).

\title[Count ladders and the boundary of the index attack]{Count ladders, duality
renormalization, and the exact boundary\\ of the index attack on the Riemann Hypothesis\\
{\normalsize\textsc{draft}}}
\author{James Scully}
\address{}
\email{jamesjscully@gmail.com}
\date{June 2026 (draft)}
\subjclass[2020]{11M26 (Primary); 19D55, 55P91, 14G40 (Secondary)}
\keywords{Riemann hypothesis, Weil positivity, topological Hochschild homology,
real cyclotomic spectra, regularized determinants, Arakelov geometry, Hodge index}

\begin{document}

\begin{abstract}
We report the outcome of a systematic attempt to transplant the function-field proof of the
Riemann Hypothesis --- intersection square, Hodge index, Riemann--Roch, effectivity, Frobenius
family --- to the integers, executed under a fail-closed discipline in which every claim is
graded (theorem / certified / census / reading) and every numerical assertion is gated by an
exit-coded witness. The supply side of the completed zeta function is reconstructed from
integer cardinalities of topological Hochschild homology with no zeta input: the von Mangoldt
atoms are the orders $\log|\THH_{2i}(\F_p)|=\Lam(p^i)$, the archimedean digamma shape is the
regularized characteristic function of the B\"okstedt torsion tower, the $p^{s/2}$ self-dual
normalization is forced by ladder-negation naturality, and a connective-$\ku$ ladder carries
the global Chebyshev function in single cardinalities. The compactified local object is
realized as an honest spectrum $\Xbar$ and its self-duality is half-proven (the Lin-block
clause is a theorem with a unique duality datum). The Weil-positivity index inequality is
posed on an arithmetic square with every ingredient constructed --- index shape (certified
Lorentzian), integrality (exact), effectivity (an elementary theorem with explicit constants),
the monoid Frobenius family, and the inequality itself --- except one: a class- and
scale-coupled archimedean unit supplying a second, independent evaluation of the
weight-against-Frobenius pairing. That one object is specified by six independent derivations,
its computed shadow has period-mean $(\gamma+\log\pi)/\log p$ at each place, and the
cheaper alternatives
are excluded by theorems: functional-equation properness is blind to the missing constant,
per-cell surplus is forbidden by an exact homogeneous functional equation, integer signature
flow provably suppresses the zeros, and the two-dimensional Castelnuovo step fails by an
exponential-versus-linear depth wall. Nothing here proves RH; the contribution is the exact
location and specification of the remaining wall, together with a no-go map and roughly
seventy fail-closed computational witnesses.
\end{abstract}

\maketitle
\tableofcontents

\section{Introduction}\label{sec:intro}

\subsection{Weil positivity and the function-field model}

Weil's criterion \cite{Weil52,Bombieri00,Yoshida92} states that the Riemann Hypothesis is
equivalent to the positivity of the explicit-formula functional on a suitable class of test
functions. For curves over finite fields the corresponding positivity is a theorem, and its
proof has a fixed anatomy. Let $C/\F_q$ be a smooth projective curve and $S = C \times C$ the
surface on which the proof lives. The argument consumes exactly five ingredients:
\begin{enumerate}
\item[(F1)] \emph{an integer-valued intersection pairing on a square} --- the surface gives
torsion classes nonzero self-intersection, which a curve cannot;
\item[(F2)] \emph{the index (Hodge--Castelnuovo) shape} --- the pairing restricted to the
relevant lattice has Lorentzian signature $(1, n)$;
\item[(F3)] \emph{Riemann--Roch with Serre duality} --- $\chi(D)$ is computed from intersection
numbers, and the dual term $h^0(K - D)$ exists;
\item[(F4)] \emph{effectivity} --- $h^0(D)$ counts actual sections, so $\chi(D) > 0$ forces a
positive representative; integers cross zero-slack thresholds where real estimates cannot;
\item[(F5)] \emph{the Frobenius family with a second evaluation} --- the graphs
$\Gamma_{F^n}$ supply a two-parameter family of classes, and the one Gram entry
$\Delta \cdot \Gamma_{F^n} = \#\mathrm{Fix}(F^n)$ has a \emph{second, independent, spectral}
evaluation (Lefschetz), which is where the zeros enter.
\end{enumerate}
This paper reports the outcome of a systematic, instrumented attempt to construct each of
these five ingredients over $\Z$, carried out in 2026 as a sequence of orchestrated research
cycles (the program record is the document set cited throughout). The attempt was governed by
a standing \emph{conservation-of-difficulty} hypothesis: every shortcut to RH relocates,
rather than removes, one irreducible obstruction, and the purpose of an honest attack is to
drive that obstruction into a corner where it can be specified exactly. That is what
happened. The contribution of this paper is not a proof of anything RH-equivalent; it is
(i) an unconditional reconstruction of the supply (prime/archimedean) side of the completed
zeta function from integer cardinalities of topological Hochschild homology, with no zeta
input; (ii) a realized compactified local object with a half-proven self-duality; (iii) the
index inequality of (F1)--(F5) posed on an arithmetic square with four and a half of the five
ingredients constructed; and (iv) a battery of no-go theorems proving that the remaining half
ingredient --- a coupled archimedean unit supplying the second evaluation of (F5) --- cannot
be replaced by anything cheaper that we, or the obvious literature mechanisms, could
formulate. The missing object ends the paper specified by six independent derivations, with
its numerical shadow computed.

\subsection{Claims grading}\label{subsec:grading}

The program operated under a fail-closed epistemic discipline, and the paper preserves it.
Every labeled statement carries one of four grades:
\begin{itemize}
\item \grade{Theorem} --- proved in the cited program record at standard granularity; where
the proof rides on a literature input, that input is named at citation grade.
\item \grade{Certified} --- verified by exact (integer/rational) or ball-arithmetic (interval)
computation with a stated finite scope; the witness script is named. A certified claim is a
theorem about finitely many instances, not an asymptotic statement.
\item \grade{Census} --- verified on a registered grid (all signs machine-resolved), not for
all parameters.
\item \grade{Reading} --- an interpretation of a computation; flagged, carrying no
load-bearing weight.
\end{itemize}
Two further disciplines matter for what follows. The \emph{degrees-of-freedom audit}
(G-INT6): every dictionary between a homotopy-theoretic object and an analytic one must count
its free modeling choices (weight grading, scale, regularization scheme, ladder start) against
its independently matched constraints; a claim with no surplus of constraints is recorded as
calibration, not derivation. The \emph{fork protocol}: nontrivial zeros of $\zeta$ may be
consumed only as post-hoc comparison data, never as construction input, and every
zero-comparison runs under a freeze protocol (inputs hashed before the zeros are loaded).
Four headline self-refutations occurred during the program (Sections \ref{sec:ladders},
\ref{sec:compact}); they are reported here \emph{as corrections}, because the paper's
epistemics include its own refutations.

\subsection{Summary of results}

Table \ref{tab:ff-status} maps the five function-field ingredients to their status over $\Z$
at the close of the program.

\begin{table}[ht]
\caption{The function-field ingredients and their status over $\Z$.}\label{tab:ff-status}
\small
\begin{tabular}{p{0.30\textwidth} p{0.42\textwidth} p{0.20\textwidth}}
\toprule
\textbf{FF ingredient} & \textbf{Status over $\Z$} & \textbf{Where} \\
\midrule
(F1) integer pairing on a square & \textsc{constructed}: the Euler/Tor pairing on
$\HZ \wedge \HZ$ is sign-bearing on torsion (the curve-level pairing vanishes identically);
diagonal self-pairing $= \zeta'(0)$ in the count reading & \S\ref{sec:square} \\[3pt]
(F2) index shape & \textsc{certified}: max-Gram lemma (exact); the geometric monoid family
$\{\psi^{p^i}\}$ is Lorentzian $(1,n)$ per weight, with an $n$-stable window &
\S\ref{subsec:maxgram} \\[3pt]
(F3) RR + duality & \textsc{half-proven}: the compactified local object $\Xbar$ exists
(theorem mod one input); its self-duality conjecture DD.4 has its Lin-block clause proven
with a unique datum; the global clause is open & \S\ref{sec:compact} \\[3pt]
(F4) effectivity & \textsc{theorem}: density-RR (KK.L1) with explicit constants, elementary
(no PNT); $\chi_{\mathrm{ren}} > 0 \Rightarrow N(D) \ge 1$ exactly & \S\ref{subsec:kk} \\[3pt]
(F5) Frobenius family + second evaluation & family \textsc{constructed} (monoid, not group:
the invertible Adams family provably carries no index); the second evaluation is \textsc{the
missing object}, demanded six times independently & \S\ref{subsec:uu},
\S\ref{sec:terminal} \\
\bottomrule
\end{tabular}
\end{table}

Along the way the program proves a sequence of boundary theorems that, we believe, are of
independent interest as a map of the wall:
\begin{itemize}
\item the count grammar (nonnegative multiplicative integer data) reaches exactly the
Selberg-class landscape and no further; the Euler product, not positivity or self-duality, is
the discriminating ingredient (the Epstein boundary, \S\ref{sec:grammar});
\item the regularized characteristic function of an $S^1$-weighted count ladder produces
every local factor of the completed zeta function, with the type of factor decided by orbit
geometry (compact orbit $\to$ Euler factor; fixed point $\to$ $\Gamma$-factor); the
self-dual $p^{s/2}$ normalization is \emph{forced} by ladder-negation naturality and is
provably not a branch convention (\S\ref{sec:ladders});
\item at $p = 2$ the real place degenerates to a complex place: a new computation
(Theorem S.1) gives $\THR(\F_2)^{tC_2}$ in closed form, exhibiting the Legendre duplication
$\Gamma_\R\Gamma_\R(+1) = \Gamma_\C$ and a rank-two Barnes $\Gamma_2$ sector
(\S\ref{subsec:p2});
\item functional-equation-level properness is structurally incapable of producing the
archimedean constant: the FE layer is exactly blind to the scale (Proposition BB.1,
\S\ref{sec:global});
\item integer signature flow sees the prime-density tower of the explicit formula and is
provably blind to the zeros: each zero enters with bounded amplitude against an $O(1)$
quantization floor (\S\ref{subsec:ii});
\item the index inequality is posed, with every term defined, on the strict arithmetic square
(\S\ref{subsec:ll}); its failure on the Frobenius-graph classes is certified and
comb-determined, with no zero correlation (\S\ref{subsec:nn});
\item on the one surviving face (the connective-$\ku$ cell comb) the inequality is proven for
every prime power and every admissible scale --- and \emph{evasion-not-surplus} is a theorem:
a uniform surplus is equivalent to a positive hole coefficient, forbidden by an exact
homogeneous functional equation (SS.1--SS.5, \S\ref{subsec:ss});
\item the two-dimensional Castelnuovo step on that face fails by an exact depth wall
(exponential Frobenius depth against linear window supply), with the precise
characterization: the von Mangoldt entry occupies Weil's fixed-point slot but possesses only
one evaluation; Weil's proof works because $\#\mathrm{Fix}$ has two (\S\ref{subsec:uu}).
\end{itemize}

\subsection{What this paper does not do}

We do not prove RH, and we do not claim that completing the missing object would prove it.
The program's own guard theorem (Proposition BB.1 with its FE-blindness corollary,
\S\ref{sec:global}) cuts against the paper's natural optimism: a construction that achieves
duality and the functional equation at every place can still be vacuous for RH, because the
FE layer determines no archimedean constant and the zeros live one layer deeper, at the
density/index level. The named failure mode --- that effectivity over the compactified square
might require an input equivalent to zero repulsion --- is stated in the charter
(\S\ref{sec:index}) and remains open. What we claim is narrower and, we believe, durable:
the supply side of $\zeta$ is counting (\S\S\ref{sec:square}--\ref{sec:ladders}); the index
attack on RH over $\Z$ has been executed to a boundary that is now \emph{exact}, in the sense
that the remaining wall is one specified object, every cheaper substitute on our map is
excluded by a theorem or a certificate, and the specification (LS1--LS6 plus the
second-evaluation demand, \S\ref{sec:terminal}) is precise enough to be falsified.

\subsection{The program record and the companion paper}

The underlying record is a sequence of twenty-one program documents (17--37 in the project's
numbering) and some forty adversarially verified worker records, produced in eight
orchestrated multi-worker cycles plus a face hunt and a stress round. The discipline ledger
at close: approximately seventy witnesses, all exit-code gated and rerun by the lead at the
landing of each cycle; six headline self-corrections propagated into the record (each
appears in this paper at the point where it bit); zero fits surviving the
degrees-of-freedom audit; and every load-bearing finite numerical claim of the program's
analytic side --- the positivity certificates T1 and T2, the tracking bound $\tau < 1$ on
all fourteen support edges, and the Davenport--Heilbronn control's certified failure ---
at rigorous interval-arithmetic grade. The analytic reduction those certificates live in
(RH as an exact finite-dimensional tracking inequality along the prime filtration, with its
own no-go map) is the subject of a companion paper; the present paper is independent of it
except where the tracking quantity $\tau$ appears as a comparison observable
(\S\ref{subsec:hh}, \S\ref{subsec:nn}).

\subsection{Notation and conventions}

$\THH$, $\THR$, $\TP$ denote topological Hochschild homology, real topological Hochschild
homology, and periodic topological cyclic homology; $\Sph$ is the sphere spectrum and
$\HZ$, $H\F_p$ Eilenberg--MacLane spectra. A \emph{count ladder} is a multiset of
eigenvalue/weight data whose entries are logarithms of finite-group cardinalities;
$\prodreg$ denotes its zeta-regularized product, $\zeta_H(s, a)$ the Hurwitz zeta function
used to define it, and the \emph{anomalous dimension} of a ladder is the value
$\zeta_{\mathrm{ladder}}(0)$ governing its scale covariance ($\prodreg (c\,\lambda_n) =
c^{\,\zeta(0)} \prodreg \lambda_n$). A \emph{regularization scheme} is the additional datum
(branch/window/orientation convention) needed to evaluate $\prodreg$ on a two-sided ladder;
the \emph{symmetric} or \emph{KMS-$\tfrac12$} scheme is the one producing the self-dual
Euler factor $p^{s/2}(1 - p^{-s})$, the name recording that $n^{-1/2}$ is the KMS state
normalization of the program's earlier realization spec. $\Lam$ is the von Mangoldt
function; the \emph{explicit formula} (EF) is used throughout only as a comparison target,
never as an input. For a $C_2$-spectrum, $(-)^{C_2}$, $(-)^{tC_2}$ denote genuine fixed
points and the Tate construction; the \emph{Lin spectrum} is $Y = (H\F_2)^{tC_2}$ with
$\pi_* Y = \F_2[t^{\pm 1}]$, $|t| = -1$ \cite{GM95}. Signatures $(n_+, n_-)$ are written
with the timelike count first; \emph{Lorentzian} means $(1, n)$.

\subsection{Relation to the literature}

The genealogy is explicit. The count-level realization of local factors follows the
Deninger--Connes flow picture \cite{Deninger91,Deninger92} and its unconditional char-$p$
realization by Hesselholt \cite{Hesselholt16}; the derived square is B\"okstedt's computation
\cite{Bokstedt} read as intersection theory; the real/reflection structure rides on real
topological Hochschild homology \cite{DMPR21,QS21} and Gorenstein duality for real spectra
\cite{GreenleesMeier17,DGI06}; the $\ku$ ladder is Lee's computation \cite{Lee22}; the
curve-level integer Riemann--Roch that the square-level program models is
Connes--Consani's \cite{CCRR22,CCRR23}, and the theta-regularized effectivity is van der
Geer--Schoof's \cite{vdGS}. The number-ring control uses Lindenstrauss--Madsen \cite{LM00};
the periodic/connective fork is Stonek \cite{Stonek18}; the cooperation-lattice input is
Kane and Baker--Richter \cite{Kane81,BakerRichter05}. The blindness theorems quantify, for
this object, the standing worry articulated around the Connes--Consani Weil-positivity
program \cite{CC20}: symmetry-layer structure alone does not discriminate $\zeta$ from its
fakes. We make that worry a theorem about our own construction and use it as a guard.

\section{The count grammar and its boundary}\label{sec:grammar}

The program's first question was the cheapest one: can the explicit-formula fingerprint of
$\zeta$ be produced from \emph{counts} --- nonnegative-integer arithmetic data --- and if so,
which counts? The instrument is a pipeline \grade{Certified; \texttt{count\_\allowbreak fingerprint}}:
any integer count table $c(n)$ with $c(1) = 1$ is converted to a generalized von Mangoldt
sequence by the exact Dirichlet-log recursion
\[
\lambda(n) \;=\; c(n)\log n \;-\; \sum_{d \mid n,\ d < n} \lambda(d)\, c(n/d),
\]
then to a correlation profile regressed against $\zeta$'s fingerprint. Calibration: the
divisor function $d(n)$ (i.e.\ $Z = \zeta^2$) returns slope $+2.0000$ at machine precision
($2.2\cdot10^{-16}$); abelian-group counts and Gaussian-integer ideal counts return residuals
supported exactly on prime powers with the predicted structures. The pipeline reports which
zeros a count structure touches.

Three grammar refinements were forced by measurement and one structural counterexample:

\begin{enumerate}
\item \emph{Self-duality} \grade{Reading, then confirmed structurally}. Count structures see
zeros at the line set by their counting growth; the critical line $\mathrm{Re}\,s = 1/2$ is
reached by \emph{self-dual} count structures ($d(n)$ sits at $1/2$ because $ab = n$ is
$a \leftrightarrow b$ symmetric). In count language the functional equation \emph{is} the
duality of the pairing.
\item \emph{Temperedness}. Self-duality alone does not put zeros on the self-dual line:
the subgroup structure of $\Z^2$ is perfectly self-dual (order-$m$ subgroups correspond to
index-$m$ subgroups via annihilators), yet its zeta function $\zeta(s)\zeta(s-1)$ has zeros
provably off its self-dual line $\mathrm{Re}\,s = 1$ --- that is the prime number theorem.
What $d(n)$ has and $c(n) = n$ lacks is temperedness: subpolynomial counts at the
square-root-normalized scale.
\item \emph{The Epstein boundary} \grade{Theorem-level, by classical examples}. Epstein zeta
functions are count structures ($r_Q(n) \ge 0$ counts lattice vectors), self-dual, tempered
--- and have off-line zeros. Hence counts $+$ self-duality $+$ temperedness do \emph{not}
force the critical line; the Euler product (multiplicative closure: ideal counts multiply,
generic quadratic-form counts do not) is the irreducible discriminating ingredient.
\end{enumerate}

The combination self-dual $+$ tempered $+$ Euler is the Selberg-class shape restricted to
cardinality coefficients: the count grammar's nonnegativity, by itself, lands exactly on the
standard conjectural landscape and does not escape it. The conclusion that redirected the
program: the index, if it exists, is not the counts' entrywise nonnegativity --- it must be a
\emph{signature theorem on a pairing matrix built from the counts}, which is precisely the
function-field structure (entrywise positivity of intersection numbers is the input; the
Hodge index is a lattice-signature statement on top).

\subsection{The Hall form is a metric; the Euler form vanishes}\label{subsec:hall}

The natural integer-valued self-dual pairing on finite arithmetic objects exists: the Hall
numbers $g^G_{H,K} = \#\{N \le G : N \cong K,\ G/N \cong H\}$ are nonnegative integers,
symmetric in $(H, K)$ by Pontryagin duality ($N \mapsto \mathrm{Ann}(N)$, $G^* \cong G$ for
finite abelian $G$). Their weighted Gram matrices are exactly computable.

\begin{proposition}[\grade{Certified; \texttt{count\_\allowbreak hall\_\allowbreak signature}}]\label{prop:hall}
The self-dual-grade Hall forms are positive definite on all computed cells
($p = 2$: $m \le 3$; $p = 3$: $m \le 2$): signature $(k, 0)$. The form is metric-like; no
index shape is present at this grade (consistent with Zelevinsky's PSH positivity for Hall
structures \cite{Zelevinsky81}).
\end{proposition}

\begin{proposition}[\grade{Theorem}]\label{prop:eulervanish}
The Euler form $\chi(H, K) = \sum_i (-1)^i \log\#\Ext^i(H, K)$ vanishes identically on finite
$\Z$-modules: $\#\Ext^1(A, B) = \#\Hom(A, B)$ for finite abelian groups (apply $\Hom(-, B)$
to $0 \to \Z \to \Z \to \Z/n \to 0$: $\#B[n] = \#(B/nB)$).
\end{proposition}

Proposition \ref{prop:eulervanish} is the categorical incarnation of ``$\mathrm{Spec}\,\Z$ is
one-dimensional, no room for signature'': torsion sheaves on a \emph{curve} always pair to
zero; they acquire nonzero self-pairing only on a \emph{surface}. The redirected question ---
compute the Euler pairing on the derived category of finite complexes over the derived tensor
square $\Z \otimes^{L}_{\Sph} \Z$, where the sphere's higher torsion supplies a second
dimension --- is the subject of the next section.

\section{The derived square}\label{sec:square}

\subsection{The diagonal is \texorpdfstring{$\THH(\Z)$}{THH(Z)}}

Let $R = \HZ \wedge \HZ$ be the derived arithmetic square (smash product over the sphere
spectrum $\Sph$) and define the Euler (Tor/Serre intersection) pairing on finite $R$-modules
\[
\langle M, N\rangle_R \;=\; \sum_i (-1)^i \log \#\,\pi_i\!\left(M \wedge_R N\right),
\]
the same pairing that vanishes identically at curve level (Proposition
\ref{prop:eulervanish}). The marquee class is the diagonal $\Delta = \HZ$ ($R$-module via
multiplication), and by definition of topological Hochschild homology
\[
\Delta \wedge_R \Delta \;=\; \HZ \wedge_{\HZ \wedge \HZ} \HZ \;=\; \THH(\Z):
\]
the derived self-intersection of the diagonal in the arithmetic square \emph{is} $\THH(\Z)$
--- the object the program sought has a standard name. The input, at citation grade, is
B\"okstedt's theorem \cite{Bokstedt}:
\[
\THH(\Z) \simeq \HZ \vee \bigvee_{k \ge 1} \Sigma^{2k-1}\, H(\Z/k), \qquad
\THH(\F_p) = \F_p[\sigma],\ |\sigma| = 2 .
\]
All positive-degree homotopy of $\THH(\Z)$ is finite --- integer cardinalities. The
$\pi_0 = \Z$ is the excess part (the diagonal meets itself in a curve, not points); by
Hochschild--Kostant--Rosenberg the positive-degree torsion is the cotangent data of $\Z$
over $\Sph$.

\begin{theorem}[The sign exists, and the value is the archimedean constant;
\grade{Theorem in the count reading; witness \texttt{prime\_\allowbreak int\_\allowbreak derived\_\allowbreak square}}]
\label{thm:zetaprime}
All torsion of $\THH(\Z)$ sits in odd degrees, so every finite group contributes with a minus
sign, and
\[
\langle \Delta, \Delta\rangle_{\mathrm{tors}}
= -\sum_{k \ge 1} \log k
\;=\; \zeta'(0) \;=\; -\tfrac12 \log 2\pi \;<\; 0
\]
under zeta regularization; equivalently the regularized order of $\THH(\Z)_{\mathrm{tors}}$
is $\sqrt{2\pi}$. The curve-level control is exactly $0$; the local control is
$\langle \Delta, \Delta\rangle_{\F_p\text{-square}} = \tfrac12 \log p > 0$.
\end{theorem}

This is the first sign-bearing, torsion-sourced pairing on an arithmetic square in the
program --- the thing Proposition \ref{prop:eulervanish} proves cannot exist at curve level.
Two honesty notes are part of the statement. \emph{Scale note}: the value is
regularization-scheme-independent but scale-dependent (anomalous dimension $-(s+\tfrac12)$:
rescaling eigenvalues by $c$ shifts the constant by $-\tfrac12\log c$); the scale-invariant
content is the $\Gamma$-factor \emph{shape}. In the count reading --- cardinalities are
dimensionless, scale canonically $1$ --- the sign is a theorem; in the operator reading it is
convention until a scheme-selecting theorem fixes it (the audit
\grade{Certified; \texttt{prime\_\allowbreak int\_\allowbreak genus\_\allowbreak audit}} decided this sharply: the sign flips at
$c = 2\pi$). \emph{Genus demotion}: the reading $2 - 2g_\Z = \zeta'(0)$, giving
$g_\Z = 1 + \tfrac14\log 2\pi > 1$, spends both the unit and the reading and is recorded as
calibration; it does not contradict Connes--Consani's genus-$0$ Riemann--Roch for
compactified $\mathrm{Spec}\,\Z$ \cite{CCRR22} because the two computations are
different invariants (their cohomological $\chi$ versus this cotangent probe), a clash
settled at invariant level: no unit reconciles them
\grade{Certified; \texttt{prime\_\allowbreak int\_\allowbreak genus\_\allowbreak audit}}.

\subsection{The supply dictionary: von Mangoldt is a THH log-cardinality}

\begin{theorem}[\grade{Theorem; exact integer identity, zero free choices;
witness \texttt{prime\_\allowbreak int\_\allowbreak thh\_\allowbreak supply}}]\label{thm:supply}
For every prime $p$ and every $i \ge 1$,
\[
\log \#\,\THH_{2i}(\F_p) \;=\; \log p \;=\; \Lam(p^i).
\]
The $s$-weighted ladder sums to the local supply $-\zeta_p'/\zeta_p(s)$ exactly; atoms at
non-prime-powers, or atoms at $p$ with weight $\ne \log p$, have no carrier: a closed point
of $\mathrm{Spec}\,\Z$ with residue field of size $n$ exists iff $n$ is prime.
\end{theorem}

The last clause is the program's strongest carrier kill: the Davenport--Heilbronn function
cannot be built in this category --- not because a computation fails but because the object
does not exist. Weighting the B\"okstedt tower by the canonical circle grading (degree
$2k-1 \leftrightarrow$ weight $k$; the Connes operator is where the $k$ in
$\Z/k = \pi_{2k-1}$ comes from) gives the diagonal's regularized characteristic function
\[
D(s) \;=\; \prodreg_{k \ge 1}\,(s + k) \;=\; \frac{\sqrt{2\pi}}{\Gamma(s+1)}, \qquad
-\frac{d}{ds}\log D(s) = \psi(s+1),
\]
a digamma --- the analytic shape of the explicit formula's archimedean supply term
\grade{Certified to $4\cdot10^{-10}$/$2\cdot10^{-6}$; same witness}. Under the G-INT6 audit
the weight choice and scale are spent modeling choices; what survives invariantly is that the
log-derivative of the diagonal's characteristic function \emph{is a digamma}.

\subsection{Properness, decomposed; the wall located}\label{subsec:properness}

The function-field proof uses properness through four separable ingredients: finiteness (PF:
coherent cohomology finite-rank, Euler characteristics honest integers), duality (PD: Serre
duality, the $h^0(K-D)$ term exists), Riemann--Roch (PR: $\chi(D)$ from intersection
numbers), and sections (PS: $h^0$ counts actual objects). The derived square is ``affine over
$\Sph$'': it has PF degree-wise but not in total (the infinite tower is why a transcendental
constant appears at all), and none of PD/PR/PS. On the function-field side the zeros enter
only through global cohomology (PR/PS), never through the local supply. The classical model
of the missing package exists: van der Geer--Schoof \cite{vdGS} define $h^0$ for compactified
$\mathrm{Spec}\,\Z$ as a theta-regularized section count and prove an exact Riemann--Roch
whose entire proof is Poisson summation. Hesselholt's theorem \cite{Hesselholt16} --- for
$X$ smooth \emph{proper} over $\F_q$, regularized determinants over $\TP$ compute the
Hasse--Weil zeta function --- has properness as hypothesis, not output. Two sharpenings
from the program's verification fleet frame everything that follows. Even in the proven
char-$p$ case, properness is used only as \emph{finiteness}: the duality/index layer
(PD/PR/PS) is never exercised and the zeros are imported from Weil's theorem --- so a
candidate $\Sph$-level duality must be shown to do more than reproduce the functional
equation, and the scenario ``$\Sph$-level Poisson achieved and still vacuous for RH'' has
no other guard (this becomes a theorem about the program's own object in
\S\ref{sec:global}). And duality alone is carrier-blind: classical Poisson summation holds
for every lattice --- Epstein passes --- so a functional-equation layer can never
discriminate. Meanwhile the curve-level existence theorem is in the literature:
Connes--Consani have proved an integer-valued Riemann--Roch with Pontryagin--Serre duality
for compactified $\mathrm{Spec}\,\Z$ over their categorical base \cite{CCRR22,CCRR23},
with a clean circularity audit (no RH, no zeros, no theta input): the missing object exists
at curve level as a theorem; the index lives on the square. The wall, in one phrase:
\emph{the program needs Poisson summation as a theorem about ring spectra; everything else
on the supply side is already counts.}

\subsection{The connective \texorpdfstring{$\ku$}{ku} ladder and the number-ring control}
\label{subsec:ku}

Two further squares calibrate the picture. For the connective square $\ku \wedge \ku$, whose
diagonal is $\THH(\ku)$ (Lee \cite{Lee22}; periodic $\KU$ is excluded: inverting Bott
rationalizes the cotangent data \cite{Stonek18}), the torsion obeys the exact closed form
\grade{Theorem; three implementations, exact integers to $w \le 3000$;
witness \texttt{prime\_\allowbreak int\_\allowbreak ku\_\allowbreak ladder}}
\[
\#\,T_{2w} = \frac{\mathrm{lcm}(1..w{+}1)}{w(w+1)}, \qquad
\log \#\,T_{2w} = \psi_{\mathrm{Cheb}}(w) + \Lam(w+1) - \log w - \log(w+1),
\]
with $\psi_{\mathrm{Cheb}}$ the Chebyshev function: where the $\HZ$ square assembles the
supply place by place, the $\ku$ square's torsion carries the \emph{global} Chebyshev
function in single cardinalities --- the fracture square internalizes the sum over places.
Its torsion is all even-degree and pairs positively (the $H^2$/surface slot), opposite to
$\HZ$; the two squares' signs are reconciled by the indexing/reading bookkeeping of the
scheme audit, not contradictory.

The standing fitted-dictionary control is the number-ring family: by Lindenstrauss--Madsen
\cite{LM00}, $\#\,\THH_{2n-1}(\mathcal{O}_K) = n^d\,|d_K|$, hence
\[
\langle \Delta, \Delta\rangle_{\mathcal{O}_K} \;=\; d\,\zeta'(0) + \tfrac12 \log |d_K|,
\]
every field distinct, the discriminant entering with exponent $\tfrac12$ --- exactly the
$|d_K|^{s/2}$ of $\xi_K$'s functional equation, with ladder multiplicity
$d = r_1 + 2r_2$ the total $\Gamma_\R$-multiplicity of $\xi_K$
\grade{Certified; \texttt{prime\_\allowbreak int\_\allowbreak genus\_\allowbreak audit}}.
Measured values:
\begin{center}
\begin{tabular}{lccccc}
\toprule
$K$ & $\Q$ & $\Q(i)$ & $\Q(\sqrt5)$ & $\Q(\sqrt{-23})$ & $\Q(\zeta_5)$ \\
\midrule
$\langle\Delta,\Delta\rangle_{\mathcal{O}_K}$ & $-0.919$ & $-1.145$ & $-1.033$ & $-0.270$
& $-1.262$ \\
\bottomrule
\end{tabular}
\end{center}
A dictionary that, fed a different ring, returns that ring's own completed-zeta datum is not
fitting; any future construction claiming to produce $\zeta$'s data must pass this bench.

\section{Ladders and schemes}\label{sec:ladders}

\subsection{The full ladder and a correction}\label{subsec:fullladder}

The Tate/duality completion of the local THH half-ladder replaces $\{s + k\}_{k \ge 1}$ by
the full Frobenius-orbit ladder $s + (2\pi i/\log p)\Z$ (the orbit is a circle of
circumference $\log p$). The regularized identity \grade{Certified;
\texttt{prime\_\allowbreak int\_\allowbreak tate\_\allowbreak ladder}, assembled from independently exhibited pieces to
$3\cdot10^{-14}$} is
\[
\prodreg_{n \in \Z}\Bigl(s + \frac{2\pi i n}{\log p}\Bigr)
\;=\; 2\sinh\!\Bigl(\frac{(\log p)\,s}{2}\Bigr)
\;=\; p^{s/2}\bigl(1 - p^{-s}\bigr)
\quad \text{\emph{under the conjugation-symmetric scheme}}.
\]
The italicized qualifier is a correction, and it reshaped the program. As first recorded, the
identity was read as ``duality completion generates the self-dual (KMS-$\tfrac12$)
normalization.'' An adversarial verification pass against the primary sources refuted the
headline: Deninger's Proposition 2.3 \cite{Deninger92} --- the identity Hesselholt's theorem
runs on --- gives $1 - \lambda q^{-s}$ on the \emph{same} full ladder under his
argument-convention scheme; the three schemes differ by exactly $p^{(s-\alpha)/2}$, the
ladder's vacuum half-sum; and Hesselholt's Theorem A \cite{Hesselholt16} is consistent only
with the prefactor-\emph{free} scheme \grade{Certified;
\texttt{prime\_\allowbreak int\_\allowbreak deninger\_\allowbreak convention}}. A regularized ladder is \emph{data plus scheme}:
the completion generates the full ladder (homotopy-theoretically grounded at finite places);
the $p^{s/2}$ prefactor is a property of the symmetric regularization scheme, which at that
point no homotopy-theoretic theorem selected. Two by-products of the same audit became
standing invariants: full ladders have anomalous dimension $0$ and half-ladders
$(1-s)/2$ --- a scale-invariant form of the place dichotomy --- and the half-ladder's
$1/(2\pi)$ scale anomaly is exactly the explicit formula's $-\tfrac12\log\pi$ supply term.

The same audit cycle surfaced two sibling tensions --- the genus clash of
\S\ref{sec:square} and the opposite torsion signs of the $\HZ$ and $\ku$ squares --- and
collapsed all three to one pre-registered question (Q-SCHEME): for each square, state the
computed invariant as (ladder data, weight dictionary, scale, scheme) and verify
cross-dictionary consistency in common units. The genus clash was decided at invariant level
(\S\ref{sec:square}); the sign question was decided as reading-relative; the scheme question
became the spine of this section. The corrective lesson is recorded as method: \emph{any
claim that a construction ``produces'' an analytic normalization must name the scheme and
exhibit the theorem that selects it.}

The structural statement that survives the correction intact:

\begin{theorem}[orbit dichotomy for count ladders; \grade{Theorem at finite places via
\cite{Hesselholt16,Deninger92}; archimedean case by the fixed-point computation below}]
\label{thm:dichotomy}
One construction --- the regularized characteristic function of an $S^1$-weighted count
ladder --- produces every Euler factor of the completed zeta function, with the factor type
decided by orbit geometry: a compact orbit of circumference $\log p$ gives the full
$\Z$-ladder of spacing $2\pi i/\log p$, hence $p^{s/2}(1 - p^{-s})$ under the symmetric
scheme; the fixed point at the archimedean place gives the degenerate half-ladder of spacing
$2$, hence $\sqrt{2\pi}/\Gamma(s/2)$, whose log-derivative on the critical line is
$\tfrac12\psi(\tfrac14 + \tfrac{it}{2})$; the rank parts ($\pi_0$ excess) are the slots of
the $s(s-1)$ pole factors. No zero-side input appears anywhere in the construction.
\end{theorem}

\subsection{Reflections: the fold lives on fixed points}\label{subsec:fold}

``TP of the square'' as a compactification vessel is structurally killed by the
Bhatt--Morrow--Scholze triple \cite{BMS18}: TP technology is $p$-complete throughout, its
Frobenius is $\varphi$-semilinear (no eigenvalues over $\Z_p$), and the archimedean fracture
corner is empty --- $\sqrt{2\pi}$ is invisible to every completion. The replacement demand is
a real/$O(2)$-equivariant refinement, and its first checkpoint passed in sharpened form
\grade{Theorem at odd $p$, via \cite{DMPR21,QS21}; witness \texttt{prime\_\allowbreak int\_\allowbreak thr\_\allowbreak fold}}:
at odd $p$ the $C_2$-\emph{Tate} construction of real THH vanishes
($(H\F_p)^{tC_2} = 0$, Quigley--Shah), while the genuine $C_2$-\emph{fixed points} carry
exactly the predicted fold,
\[
\pi_*\,\THR(\F_p)^{C_2} = \F_p[y], \quad |y| = 4, \quad \res(y) = x^2
\qquad \text{(Dotto--Moi--Patchkoria--Reeh \cite{DMPR21})},
\]
i.e.\ weights $\{n \ge 0\} \to \{2n\}$: half-line, start $0$, spacing doubled --- the
$\Gamma_\R$ ladder, with anomalous dimension $(1-s)/2$, on the fixed-point carrier. The
functorial dichotomy --- rotation subgroups $\to$ Tate $\to$ full ladders (local factors);
reflection subgroups $\to$ fixed points, Tate vanishing $\to$ folded half-ladder
($\Gamma$-factor) --- is the structural form of the program's repeated finding that duality
completion does not happen at the archimedean place. The $C_2$-involution acts on
$\pi_{2k}\THH(\F_p)$ by $(-1)^k$: Serre's classical $F_\infty$ eigenvalue selection
materializes as genuine equivariance.

\subsection{Scheme rigidity: KMS-\texorpdfstring{$\tfrac12$}{1/2} is forced by negation,
and by nothing else}\label{subsec:rigidity}

The reopened scheme question closed in two steps, the first of which corrected the second's
own first formulation. The hypothesis ``the $C_2$ of THR acts on the determinant layer''
is satisfiable-but-vacuous: the residual Weyl action satisfies it for \emph{every} scheme.
The forcing form (H-r3) demands that the $C_2$ act by ladder \emph{negation}
$N : (V, \Theta) \mapsto (V, -\Theta)$ covering $s \mapsto -s$, with the regularized
determinant natural for it. The analytic half of that hypothesis is a theorem
\grade{Theorem; witness \texttt{prime\_\allowbreak int\_\allowbreak hr3\_\allowbreak scheme\_\allowbreak rigidity}}:

\begin{lemma}[parity lemma]\label{lem:parity}
In the vacuum-term scheme family $f_c(w) = q^{cw}(1 - q^{-w})$:
$f_c(-w) = -\,q^{(1-2c)w} f_c(w)$. Hence $c = \tfrac12$ --- the symmetric/KMS-$\tfrac12$
scheme --- is the unique scheme equivariant under the reflection $w \mapsto -w$.
\end{lemma}

\begin{proof}
$f_c(-w) = q^{-cw}(1 - q^{w}) = -\,q^{-cw} q^{w} (1 - q^{-w}) = -\,q^{(1-2c)w} f_c(w)$;
the prefactor is $w$-independent iff $c = \tfrac12$.
\end{proof}

\begin{theorem}[scheme rigidity, A0--A2]\label{thm:rigidity}
\textup{(A0)} Every eventually-windowed branch scheme on the full ladder has value
$D_{a_+,a_-}(w) = (-1)^{a_+ + a_-}\, q^{-(a_+ - a_-)w}\,(1 - q^{-w})$, $(a_+, a_-) \in \Z^2$:
all branch-convention vacuum slopes are \emph{integers}.
\textup{(A1)} If $f(w) = e^{P(w)}(1 - q^{-w})$ satisfies negation-naturality
$f(-w) = u f(w)$ for a constant unit $u$, then $P_{\mathrm{odd}}(w) = \tfrac{\log q}{2} w$
exactly and $u = -1$: every negation-natural scheme value is
$e^{a}\, 2\sinh(w \log q / 2)$, the symmetric scheme up to an EF-invisible constant.
\textup{(A2)} Consequently no pointwise branch convention --- Deninger's, Hesselholt's, or
any other --- can ever produce the KMS-$\tfrac12$ prefactor: naturality demands slope
$\tfrac12$ and conventions supply integers. The symmetric scheme is not a branch choice at
all.
\end{theorem}

\begin{proof}[Proof of (A1)]
$1 - q^{w} = -q^{w}(1 - q^{-w})$, so $f(-w) = u f(w)$ reads
$e^{P(-w) - P(w) + w\log q + i\pi} = u$, i.e.\ $-2P_{\mathrm{odd}}(w) + w\log q$ is
constant. A polynomial identity: $P_{\mathrm{odd}}(w) = \tfrac{\log q}{2}w$, the constant is
$0$, and $u = e^{i\pi} = -1$. (Why $u$ must be constant: in finite dimensions the unit is
$(-1)^{\dim V}$; allowing $u = u(w)$ would permit every scheme and is naturality of
nothing.) For (A0), split the vertical ladder at its two ends and evaluate the windowed
Hurwitz regularizations; the Gamma reflection identity assembles the stated closed form ---
the full computation is in the record, witnessed at $3\cdot10^{-15}$. (A2) is immediate
from (A0) $+$ (A1).
\end{proof}

Three companion propositions complete the audit \grade{Theorem; same witness}:

\begin{proposition}[A3--A6]\label{prop:a3a6}
\textup{(A3)} Deninger's Convention (2.1) of \cite{Deninger92} is a circle-orientation
choice (a generator of the deck group of the log cover); under negation the two ladder ends
swap and the window maps to a different convention --- no single window is negation-stable.
The $C_2$-average of Deninger's scheme and its mirror is the symmetric scheme, with the
$\mu_4$ sign-count ambiguity resolved to $+1$ by the pair-modulus construction.
\textup{(A4)} Conjugation selects nothing: every real-$c$ scheme respects the Schwarz
pairing. (This corrects the program's earlier gloss: the selecting involution is negation,
the local functional-equation flip, not conjugation.)
\textup{(A6)} The symmetric determinant requires a canonical $\lambda^{1/2}$: it lives on a
double cover of the eigenvalue space, so any carrier of the forcing structure must be
2-primary.
\end{proposition}

The open residue, precisely posed, was the homotopy half: does the geometric $C_2$ actually
supply negation to the determinant layer? Its resolution --- negative at THR itself,
positive at the compactified object --- is Section \ref{sec:compact}.

\subsection{The \texorpdfstring{$p = 2$}{p=2} anomaly, computed}\label{subsec:p2}

At $p = 2$ the fixed points do not fold cleanly, and the gap was discharged by a new
computation:

\begin{theorem}[S.1; \grade{Theorem additively; ring form up to a stated generator-choice
caveat; witness \texttt{prime\_\allowbreak int\_\allowbreak thr\_\allowbreak p2\_\allowbreak anomaly}}]\label{thm:s1}
$\THR(\F_2)^{tC_2} \simeq \prod_{n \ge 0} \Sigma^{2n} (H\F_2)^{tC_2}$, with
$\pi_* = \F_2[[z]][t^{\pm 1}]$ \ ($z = t^2 x$ of weight $1$, $|z| = 0$). The mechanism:
the DMPR genuine splitting of $\THR(\F_2)$ forces Borel formality, so the Tate spectral
sequence collapses.
\end{theorem}

The $p = 2$ fixed points factorize exactly as (fold core $\F_2[\bar{x}^2]$, the odd-$p$
$\Gamma_\R$ ladder) $\times$ (a $\{1, \bar x\}$ doubling) $\times$ (the Euler tower
$\F_2[y]$), and each factor is an identification: the doubling is the mod-2 degeneration of
Serre's $(-1)^n$ selection, with the exact regularized identity
$\prodreg(s+n) = \prodreg(s+2n)\cdot\prodreg(s+2n+1)$ mirroring Legendre duplication
$\Gamma_\R(s)\Gamma_\R(s+1) = \Gamma_\C(s)$ --- \emph{at $2$ the real place degenerates to a
complex place}, a computed statement, not an analogy. The Euler tower has no odd-$p$
analogue: the full degree carrier is the Barnes double Gamma $\Gamma_2(s \mid 1, 2)$ with
quadratic anomalous dimension $\zeta_B(0, s \mid 1,2) = s^2/4 - 3s/4 + 11/24$ --- the first
rank-two archimedean object the program produced from homotopy theory, the natural candidate
slot for the rank-two $s(s-1)$ pole pair. Moreover, at $p = 2$ --- the unique prime with
nonvanishing reflection Tate --- duality completion \emph{does} happen at the reflection,
and per weight $n$ the full line gives $2\sinh((s+n)\log 2/2) = 2^{(s+n)/2}(1 - 2^{-(s+n)})$:
the self-dual local factor appearing geometrically. (A later audit weakened this honestly:
the form is scheme-\emph{applied}, not scheme-forced, at this exhibit; the forcing came only
with the glued object of \S\ref{sec:compact}.) Three independent results --- the
$\lambda^{1/2}$ double cover, the geometric KMS-$\tfrac12$ at the reflection sector of $2$,
and Connes--Consani's canonical divisor $K = -2\{2\}$ supported at the prime $2$
\cite{CCRR22} --- point at the same address.

\subsection{The index shape: the max-Gram lemma and the certified Lorentzian family}
\label{subsec:maxgram}

The Frobenius-family search produced a clean negative, its explanation, and the repair. The
bare Adams family $\{\psi^1, \dots, \psi^n\}$ on the $\ku$ square has per-weight pairing
cardinalities $|k - l|$, and its log-Gram is \emph{not} Lorentzian at any $n \ge 4$
\grade{Certified, refutation-grade; \texttt{prime\_\allowbreak int\_\allowbreak adams\_\allowbreak family\_\allowbreak signature}}. The
explanation is a theorem within the $p$-complete shadow: all invertible Adams graphs have
bidegree $(1,1)$ --- the family is a group, every pairing cardinality is a local Lefschetz
factor $\det(1 - d\psi^\alpha \mid \mathrm{wt}\ w) = p\text{-part of } \alpha^w - 1$ --- so
the Castelnuovo degree-growth mechanism provably cannot come from the invertible family; it
requires the non-invertible $\psi^k$, i.e.\ the multiplicative \emph{monoid}. The repair:

\begin{lemma}[max-Gram lemma; \grade{Theorem}]\label{lem:maxgram}
Let $a_0 < a_1 < \cdots < a_n$ be real with $a_n > 0$, and let
$M[i,j] = \max(a_i, a_j)$, $\MIN[i,j] = \min(a_i, a_j)$, $u = (a_i)$. Then $M$ has
signature $(1, n)$.
\end{lemma}

\begin{proof}
$\max(a_i,a_j) + \min(a_i,a_j) = a_i + a_j$ gives
$M = u\mathbf{1}^{T} + \mathbf{1}u^{T} - \MIN$. The rank-two part
$u\mathbf{1}^T + \mathbf{1}u^T$ has signature $(1,1)$, and $\MIN$ is positive semidefinite
(positive definite for positive $a_i$: $\det\MIN = a_0\prod(a_{i+1} - a_i)$, the nested-min
factorization), so $n_+(M) \le 1$ by Weyl subadditivity. Since $M[n,n] = a_n > 0$,
$n_+ \ge 1$; and $\det M = (-1)^n a_n \prod(a_{i+1} - a_i) \ne 0$ kills the kernel.
\end{proof}

The Castelnuovo/Hodge-index shape is therefore \emph{automatic} for any family whose
log-pairing is $\max(\deg_i, \deg_j)$ --- and the arithmetic family carrying that shape is
the geometric one $\{\psi^{p^0}, \dots, \psi^{p^n}\}$:
$\log|p^{iw} - p^{jw}| = (w\log p)\max(i,j) + E$ with $\|E\|$ exactly bounded
($\|E\|_\infty \le 2x/(1-x)^2$, $x = p^{-w}$), certified Lorentzian per weight for all
$(p, w)$ with $w \ge 3, 2, 2, 1$ at $p = 2, 3, 5, \ge 7$, with an $n$-stable window of
admissible diagonal shifts containing $0$
\grade{Certified; \texttt{prime\_\allowbreak monoid\_\allowbreak frobenius\_\allowbreak shadow}}. The same pairing read
$p$-adically is the $\MIN$/ultrametric PSD face: by the product formula
$|p^{iw} - p^{jw}| = p^{w\min(i,j)}(p^{w|i-j|} - 1)$, one pairing has two faces ---
$v_p$-Gram $= w\MIN$ (hierarchical, PSD) and log-Gram $= (w\log p)\MAX + E$ (the index
universe) --- split by the product formula. Groups carry no index shape because they carry
no degrees. Separately, the branch fork on the twisted family was closed from the
literature: by Kane (19.6.1) \cite{Kane81} via Baker--Richter \cite{BakerRichter05} the
cooperation lattice is monomial through $w \le p^2 - 2$, so the local Euler supply does
\emph{not} live in the invertible twisted Adams family.

\section{The compactified object}\label{sec:compact}

\subsection{The product-formula determinant spec}

The program's repeated lesson was that imposing a known formalism (TP, branch conventions,
tolerance categories) yields the right ladders in the wrong layer. The two structures that
never failed a kill-test are the count identities and the product formula; the compactified
object was therefore \emph{posed} as the structure they define (the analogy to how cyclotomic
spectra were defined --- Frobenius as structure, existence as theorem --- is deliberate). A
\emph{product-formula determinant} (PFD) object consists of: the full determinant ladders
$(V_p, \Theta_p)$ of the $p$-local Tate squares; the degenerate reflection sector
$(V_\infty, \Theta_\infty)$ (half-ladder, spacing $2$; at $2$ the computed $\Gamma_\C$
degeneration and Barnes $\Gamma_2$ tower of Theorem \ref{thm:s1}); the product formula
imposed as the gluing identification ($\prod_v |\lambda|_v = 1$: the $p$-adic MIN faces and
the archimedean MAX face declared two readings of one object); and Poisson as defining
property --- the finite-truncation residue pairings must \emph{extend}. Pre-registered kill
tests: K1 (the $1/(2\pi)$ scale must emerge, never be inserted), K2 (the duality must restore
$s \mapsto -s$ exactly where completion breaks it, which by Theorem \ref{thm:rigidity} forces
KMS-$\tfrac12$), K3 (the Barnes $\Gamma_2$/$s(s-1)$ rank sector must appear), K4 (the
quantization-defect pole must be evaded), K5 (carrier: no ring, no object; the
$\mathcal{O}_K$ bench must return $\xi_K$'s data), K6 (zero free gluing choices).

Two prior negatives shaped the spec, and both are part of the record. First, the F1 probe at
THR returned \emph{negative}: the residual Weyl action \emph{shears} (Artin--Schreier-style),
it never negates; the inverting element is consumed by the Tate construction; so realization
input KMS-$\tfrac12$ has no carrier at THR itself
\grade{Certified; \texttt{prime\_\allowbreak int\_\allowbreak u\_\allowbreak f1\_\allowbreak probes}}. What survived: negation-as-duality
exists at every finite weight truncation (a perfect residue pairing $m \leftrightarrow N-m$)
and \emph{dies in the completion} --- Poisson-shaped to its core. Second, the original
``equivariant determinant layer'' hypothesis was vacuous as first worded
(\S\ref{subsec:rigidity}); the forcing form replaced it.

\subsection{The gluing, executed; the scheme forced}

\begin{theorem}[the $p=2$ instance; \grade{Theorem at the $\pi_*$/filtered level, riding on
Theorem \ref{thm:s1}; witness \texttt{prime\_\allowbreak int\_\allowbreak z\_\allowbreak u2\_\allowbreak glue}}]\label{thm:glue}
Let $D$ be the weight-completed disk of Theorem \ref{thm:s1} and $D'$ its truncation tower's
Matlis dual ($D'$ is \emph{forced}: the unique compatible system extending the finite-level
residue pairings). The glued object $G = D \cup_{\mathrm{Lin}} D'$ exists with zero free
choices; compactification restores negation as a duality; Theorem \ref{thm:rigidity}(A1)
applies as a structure map; and $c = \tfrac12$ \emph{is forced} at $p = 2$. The shear of the
F1 probe is identified as $[-1]$ of the multiplicative formal group, giving the unit
$u = -1$ of (A1) a structural carrier.
\end{theorem}

The gluing is not 2-specific: executed at $p = 3, 5$ with $c = \tfrac12$ forced and the
structural deltas being absences (no Euler tower)
\grade{Certified; \texttt{prime\_\allowbreak int\_\allowbreak gc\_\allowbreak oddp\_\allowbreak glue}}. The K1 gate resolved as
\emph{not-at-one-place} with the mechanism exact: the glued object is scale-free
($\zeta_H(0,a) + \zeta_H(0,1-a) = 0$; the glued trace is a lattice delta at $2^{-s} = 1$, the
local face of the product formula), so the $1/(2\pi)$ is a cross-place effect --- which is
what the product-formula gluing is for. K3 passed (the Barnes sector enters $G$, with
$\zeta_2(0, s+1 \mid 1,2) = s(s-1)/4 - 1/24$ self-dual under $s \mapsto 1-s$ at exactly $G$'s
forced recentering); K4 passed by a mechanism neither registered reading anticipated:
$G$'s atoms are exactly the pole-support set of the quantization defect (\S\ref{sec:global}),
each branch carries the pole, and the K2 duality \emph{cancels the pole pair exactly}
(residue $0$ to $10^{-86}$) --- properness as pole-cancellation
\grade{Certified; \texttt{prime\_\allowbreak int\_\allowbreak gc\_\allowbreak k4\_\allowbreak integer\_\allowbreak layer}}.

\subsection{The strict model}

The PFD structure has a model with zero free choices inside the strict (non-derived) layer
\grade{Theorem; witness \texttt{prime\_\allowbreak cc\_\allowbreak thh\_\allowbreak norm\_\allowbreak filtration}}:
transporting Connes--Consani's total-mass filtration along the strict square's universal
character (whose target is the monoid ring of $\Z\setminus 0$ --- the refutation-grade fact
that the strict square's \emph{total} memory is the multiplicative monoid, the program's
third independent arrival at the Euler-product carrier) makes the strict cyclic bar a
norm-filtered cyclic object. The norm \emph{is} the MAX face:
$N([p^i] - [p^j]) = 2\,p^{\max(i,j)}$, its Gram on the geometric family is certified
Lorentzian with no low-weight anomaly, the $v_p$ of the same classes are the MIN faces, and
$\log N = \sum_p v_p \log p$ exactly --- the product formula as an identity of invariants.
Epstein keeps exactly the carrier-blind half (its Lorentzian \emph{pass} is the lesson:
index shape alone discriminates nothing) and dies at $\mu$, $v_p$, and the product formula.

\subsection{The spectrum exists; the duality half-proven}

\begin{theorem}[DD.1--DD.3; \grade{Theorem mod Theorem \ref{thm:s1}; witness
\texttt{prime\_\allowbreak int\_\allowbreak dd\_\allowbreak bc\_\allowbreak dual}}]\label{thm:dd}
\textup{(DD.1)} No honest spectrum $T$ with $\pi_* T = G$ satisfies a plain Brown--Comenetz
self-duality $D_{\F_2}(T) \simeq \Sigma T$: the full dual of $\F_2[[z]]$ is too big
(Erd\H{o}s--Kaplansky).
\textup{(DD.2)} The continuous Brown--Comenetz dual of the $z$-adic truncation tower matches
the forced $D'$ exactly: the dual tower is isomorphic to the multiplication-by-$z$ tower via
the level-$N$ residue self-dualities, with the divided-power multiplication falling out of
the dual Hopf structure --- the ``divided-power dual'' is a theorem at the filtered level.
\textup{(DD.3)} $\Xbar := \fib\bigl(X \times X' \to Y\bigr)$, with
$X = \THR(\F_2)^{tC_2}$, $Y = (H\F_2)^{tC_2}$ the Lin spectrum, and
$X' = \colim$ of the dual truncations, is an honest spectrum with $\pi_* \Xbar = G$ exactly
(Mayer--Vietoris; zero free choices, $\Aut_Y(Y) = \F_2^\times = 1$).
\end{theorem}

\begin{conjecture}[DD.4, continuous self-duality; display as corrected]\label{conj:dd4}
$\Xbar \simeq \Sigma\, I^{\mathrm{cont}}(\Xbar)$ in the appropriate filtered/pro-ind
category, the equivalence exchanging the two branches and covering $s \mapsto -s$. (The
complete/torsion home is Dwyer--Greenlees--Iyengar \cite{DGI06}; the nearest theorem-grade
precedent is $C_2$-equivariant Anderson self-duality for real spectra
\cite{GreenleesMeier17}.)
\end{conjecture}

\begin{theorem}[EE.1, the Lin-block clause of DD.4; \grade{Theorem; witness
\texttt{prime\_\allowbreak int\_\allowbreak ee\_\allowbreak dd4\_\allowbreak lin}}]\label{thm:ee1}
On the Lin block $Y = (H\F_2)^{tC_2}$ the duality datum exists, is induced by the
$t^{-1}$-coefficient via Brown--Comenetz representability, and is the \emph{unique} nonzero
$Y$-module map $Y \to \Sigma I(Y)$ --- a singleton, torsor-free as a theorem. It induces
exactly the $m = 0$ component of the K2 pairing ($(j, 0) \mapsto (1-j, 0)$); the dualizing
character is $(\mathrm{degree}, \mathrm{weight}) = (1, 0)$, the residue line.
\end{theorem}

The same worker's executed Borel computation refuted, fail-closed, the registered
identification of the $t$-shadow (Lemma N1-b as posed): the relevant bigraded slot
$\pi^{C_2}_{(-2,-1)}\THR(\F_2)$ vanishes, so any composite that forgets the parametrized
circle kills the comparison class by naturality (Theorem EE.2); the identification question
was re-posed on the geometric route with a torsor-invariant leading-term form. The duality's
remaining open content is the weights-$\ge 1$ branch-exchange clause; the torsor of
ambiguities is supported there and is one object across its three appearances.

\subsection{The scale is expelled from torsion}

\begin{proposition}[FF.1; \grade{Theorem, conditional on Conjecture \ref{conj:dd4}; witness
\texttt{prime\_\allowbreak int\_\allowbreak ff\_\allowbreak kappa\_\allowbreak xbar}}]\label{prop:ff1}
The Matlis/Brown--Comenetz exchange is the torsion shadow of the Pontryagin pair and forgets
exactly the scale carrier: the continuum is BC-annihilated, $2\pi$-rigidity is a topological
datum ($\Aut_{\mathrm{top}}(\R/\Z) = \{\pm 1\}$) with no torsion shadow. Hence no torsion-level
duality, including DD.4, can determine the archimedean scale; the scale demands a
rational/connected (Anderson-layer) carrier.
\end{proposition}

This drops DD.4's duty roster from triple (scheme, pole, scale) to double (scheme, pole) and
sends the scale to the global assembly --- where the next section proves it cannot be read
from the functional-equation layer at all.

\subsection{State at the closure of the line}

What the compactification line proved and realized: the local object $\Xbar$ exists
($\pi_* = G$, modulo Theorem \ref{thm:s1}); KMS-$\tfrac12$ is forced at $p = 2, 3, 5$; the
duality's Lin clause is a theorem with a singleton datum; the pole cancellation is
place-uniform and globally load-bearing; and every loose constant was adjudicated into a
scheme slot or a vacuum term (the $\tfrac14$-versus-$\tfrac12$ recentering is the
$\Gamma_\R/\Gamma_\C$ boundary; the $-\tfrac1{24}$ is the $\zeta(-1)$ vacuum,
center-value exact). What remained open, named: (A) the weights-$\ge 1$ clause of
Conjecture \ref{conj:dd4}, carrying the scheme and pole duties --- homotopy-theoretic,
sharply posed, one clause already a theorem; and (B) the archimedean unit, at that point
\emph{triply} specified (read the scale; supply one continuum unit per place; carry the
rational layer torsion provably lacks). The index attack of \S\ref{sec:index} was chartered
on top of this state, and added three more independent specifications of (B).

\section{The global assembly and the blindness theorems}\label{sec:global}

\subsection{The equivariance no-go and FE-blindness}

Assemble the global determinant from the banked inputs: the per-place full-ladder layers
$p^{w/2}(1 - p^{-w})$, the degenerate archimedean sector at scale $\kappa$, and the
functional-equation pairing. The gate observable is the constant
$c_\infty(\kappa) = -\tfrac12\log\kappa$ multiplying the degenerate sector; the explicit
formula requires $\kappa = \pi$.

\begin{proposition}[BB.1, equivariance no-go; \grade{Theorem}]\label{prop:bb1}
Witness: \texttt{prime\_\allowbreak int\_\allowbreak bb\_\allowbreak global\_\allowbreak assembly}.
Every banked input of the global assembly is invariant under the rescaling
$R_\lambda : \kappa \mapsto \lambda\kappa$: the full-ladder anomaly is $0$, the glued trace
identity $\frac{1}{1-x} + \frac{1}{1-1/x} - 1 = 0$ holds in exact rational arithmetic, the
product formula is integer-exact and contains only $\{\log p\}$, and the residue pairings are
$\F_2$-exact. The observable transforms as
\[ D_{\infty,\lambda\kappa}/D_{\infty,\kappa} = \lambda^{-(1-s)/2}. \]
An observable that moves
under a group fixing every input is not a function of the inputs: \emph{no constant is
determined}.
\end{proposition}

\begin{corollary}[FE-blindness]\label{cor:feblind}
The paired archimedean anomaly satisfies
$\zeta_H(0, s/2) + \zeta_H(0, (1-s)/2) = \tfrac12$ exactly, so $\kappa$ enters the
FE-paired factor only as an $s$-independent constant: \emph{the functional equation of the
assembly holds for every $\kappa$}. A global duality that only reproduces the functional
equation is structurally incapable of producing the archimedean constant; the constant lives
one layer deeper, at the critical-line density (index/zero-counting) layer, where
$-\tfrac12\log\kappa$ is visible.
\end{corollary}

Corollary \ref{cor:feblind} is the quantitative form of the program's standing vacuity guard
(classical Poisson holds for every lattice --- Epstein passes --- so a functional-equation
layer alone can never discriminate), and it is the caveat this paper applies to itself in
\S\ref{sec:terminal}. The missing ingredient acquired its first name here: the
\emph{archimedean Poisson exchange}, the comb $\leftrightarrow$ continuum duality whose
classical shadow is the Gaussian's self-duality in Tate's thesis. The $2\pi$ is \emph{in}
the object --- the Abel-regularized glued trace of each $G_p$ is exactly the comb
$(2\pi/\log p)\sum_n \delta_{2\pi n/\log p}$, Poisson summation on $(\log p)\Z$
\grade{Certified} --- but uncoupled.

\subsection{The quantization-defect pole and its cancellation}

A strict (integer-dimension) tolerance layer cannot hold the archimedean constant for an
exact reason \grade{Theorem; witness \texttt{prime\_\allowbreak cc\_\allowbreak quantization\_\allowbreak defect}}: the
quantization defect $\sum_k (\lceil \log_b k\rceil - \log_b k)\, k^{-s}$ has a simple pole at
$s = 0$ with residue $-1/(2\log b)$, carried exactly by the atoms $\{b^j\}$ --- so
integer-dimension quantization can never hold $\zeta'(0)$. The glued object's atoms are
exactly this pole-support set, and the K2 duality cancels the pole pair branch against
branch (\S\ref{sec:compact}): properness as pole-cancellation, the program's third
independent mechanism at the same slot.

\subsection{The cancellation globalizes; the zero mode}

The duality-cancellation mechanism is place-blind and is the \emph{existence} mechanism for
the global trace: uncancelled, the assembly is non-summable
\grade{Certified; \texttt{prime\_\allowbreak int\_\allowbreak gg\_\allowbreak global\_\allowbreak cancellation}}. The surviving layer
carries an exact zero-mode obstruction: \emph{one continuum unit per place} must be
subtracted, after which the renormalized survivor is the $\sigma = 0$ shell of the explicit
formula's prime side (matched to $7\cdot10^{-85}$). Duality-as-existence and the
per-place continuum subtraction are one demand seen at two layers; the subtraction is the
second independent specification of the missing unit.

\subsection{Integer signature flow is blind to the zeros}\label{subsec:ii}

The cheapest conceivable discreteness-sees-spectrum claim --- that integer signature jumps of
the growing all-places norm-Gram family correlate with zero positions --- was tested under a
freeze protocol and refuted, with the mechanism derived
\grade{Census $+$ derived mechanism; witness \texttt{prime\_\allowbreak int\_\allowbreak ii\_\allowbreak signature\_\allowbreak flow}}:
\begin{itemize}
\item the global glued family keeps the Lorentzian shape ($n_+ = 1$ at every flow sample up
to $N = 3321$, ball-certified at five instances);
\item the dimension-normalized jump-loci statistics \emph{separate} the $\mathrm{li}(x)$
density model from the crude $x/\log x$ model (separation $2.6$ counts) --- the flow sees the
smooth density tower of the explicit formula;
\item against the first twelve zeros the decision channel returns $z \le 1.7$ (threshold
$4$): no hits; the loci accumulate at $\log 2$, a norm constant, not a zero position;
\item the mechanism: each zero enters the counting function with bounded amplitude
($\sim 0.02$ counts) against the $O(1)$ quantization floor of an integer observable, while
the density tower grows; more range cannot help.
\end{itemize}
The design consequence is a theorem-shaped constraint: an integer observable that touches the
zeros must have \emph{growing} zero-response --- a section count $h^0$, not a signature. This
is the bridge to the index attack.

\section{The index attack, executed to its boundary}\label{sec:index}

The charter, fixed in advance: the conserved difficulty is \emph{effectivity over the
compactified square} (integer Riemann--Roch at square level); the index shape, integrality,
and the Frobenius family are held from the preceding sections; the Epstein lesson governs
(signature alone is carrier-blind; the discriminating conjunction is index $+$ integrality
$+$ effectivity $+$ monoid family); and the named failure mode is recorded before execution:
\emph{if effectivity on the square requires an input equivalent to zero-repulsion, the
conservation law wins.}

\subsection{The Weil form factors through the index face}\label{subsec:hh}

\begin{theorem}[HH.1; \grade{Theorem, exact rational arithmetic; witness
\texttt{prime\_\allowbreak hh\_\allowbreak weil\_\allowbreak index\_\allowbreak factorization}}]\label{thm:hh1}
In the coefficient space of the geometric family $\{[p^i]\}_{i = 0..n}$, the difference
classes $d_k = [p^k] - [p^{k-1}]$ diagonalize the certified index face:
$\langle d_k, d_l\rangle_{(w\log p)\MAX} = -(w \log p)\,\delta_{kl}$. The unit class $[1]$
is null (lightlike); $\mathrm{span}\{d_k\}$ is the degree-zero hyperplane, maximal spacelike;
its orthocomplement is timelike. In particular, at $w = 1$,
\[
\Lam(p^k) \;=\; \log p \;=\; -\,\langle d_k, d_k\rangle_M \quad \text{for every } k:
\]
the von Mangoldt atom weight is minus the self-intersection of the ladder step.
\end{theorem}

\begin{corollary}[the prime side factors exactly]\label{cor:primeside}
With the sample map $T_p(g) = \sum_k \check g(k \log p)\, d_k$ and the dual vector
$\Theta_p = 2\sum_k p^{-k/2} d_k$ (a pure KMS-$\tfrac12$ geometric series --- zeta is the
\emph{untuned} case: no per-atom retuning, the entire weight system $\Lam$ absorbed by the
Gram), the explicit-formula prime side factors exactly through the index face:
\[
P_R(g) = -\sum_p \bigl\langle T_p(g), \Theta_p\bigr\rangle_{(\log p)\MAX}
\quad \text{(defect zero)}, \qquad
Q(g) = W_\infty(g) + \sum_p \bigl\langle T_p(g), \Theta_p\bigr\rangle_M .
\]
The defect of
the index-face factorization is the archimedean supply $W_\infty$, \emph{identically}.
\end{corollary}

The geometry this exposes, each clause exact or census-grade as marked: the image of $T$
avoids the positive cone entirely (the Perron direction --- the thing the Castelnuovo
mechanism is built around --- is idle in the Weil form); the image cone is the moment cone
(sample sequences are PSD-Toeplitz, the Connes--Consani Toeplitz mechanism \cite{CC20} now
carrying an $M$-geometry); on each support rung, RH-at-the-rung is a \emph{dual-cone
membership}: $(1, \Theta)$ lies in the dual of the image cone
$K_R = \{(W_\infty(g), T(g))\}$, with slack equal to the program's certified tracking
quantity $\tau < 1$. The stronger ``positive archimedean budget'' reading \emph{fails},
measured both ways \grade{Census}: the $W_\infty$ block goes indefinite from support
$R = 2.30$ and falls while joint margins stay positive --- Weil positivity is a genuine
\emph{joint cancellation}, so the archimedean unit must couple, not patch. The
self-intersection identity also kills Davenport--Heilbronn twice over: negative atom weights
violate $\Lam = -\langle d, d\rangle \ge 0$, and non-prime-power atoms have no tower address.
And the $s = 1$ pole appears as geometry: $\sum_p \|\Theta_p\|^2_{-M}$ diverges like
$-4\zeta'/\zeta$ near $1$ --- the KMS vector is not in the lattice completion.

\subsection{Effectivity is a theorem}\label{subsec:kk}

On the strict square (the free effective monoid $M = (\Z_{\ge 1}, \times)$ with
$v_p$-coupled twist classes $D = (a; m)$, theta-regularized count
$h^0(a; m) = \log(1 + 2\sum_A e^{-\pi w_D(A)^2})$):

\begin{theorem}[KK.L1, density-RR;
\grade{Theorem with explicit constants; witness \texttt{prime\_\allowbreak int\_\allowbreak kk\_\allowbreak l1\_\allowbreak density\_\allowbreak rr}}]
\label{thm:kk}
\textup{(A)} \emph{Effective cone $=$ finiteness locus}: $h^0(a; m) < \infty$ iff
$m_p \le 0$ for all $p$; for any $m$ with some $m_p \ge 1$ the count diverges for every
$a$. Genuinely square divisors have no internal RR dual; the unique duality-stable slice is
$m = 0$, the curve shadow, where the van der Geer--Schoof RR \cite{vdGS} holds exactly.
\textup{(B)} For $m \le 0$ with twist support of size $J$:
$\chi_{\mathrm{tw}}(a; m) - a \to \log\rho(m) = \sum_p \log\rho_p$ with defect
$O(a^{J-1}e^{-a})$ and all constants explicit --- the RR constant is \emph{additive over
places}: ``one continuum unit per place'' is a theorem on the strict square's count layer.
\textup{(C)} $|N_{S_m}(x) - \rho x| \le E_T(x)$ with $E_T$ explicit. Consequently
$\chi_{\mathrm{ren}} > 0 \Rightarrow N(D) \ge 1$ exactly: effectivity fires at integer
grade, and $\log N(D)$ is the growing integer observable demanded by
\S\ref{subsec:ii}.
\end{theorem}

The proof is elementary in the strict sense: unique factorization at the twisted places,
inclusion-exclusion, second-order Euler--Maclaurin, geometric series --- no PNT, no Mertens,
no Poisson, no zeta input. The place where PNT-strength input would become necessary is
located exactly (infinite twist support). Three companion statements sharpen the boundary:

\begin{proposition}[KK.L1-R, the real wall]\label{prop:kkr}
Extending $m$ to real values, $h^0(a; m) < \infty$ iff $m_p < 1$ for all $p$; the divergence
locus is the closed set $\{\max_p m_p \ge 1\}$, whose boundary wall carries the blow-up of
the RR constant ($\rho_p \to +\infty$ as $m_p \to 1^-$). The first integer class outside the
effective cone in any positive direction sits exactly on the wall, and the in-family duals
of the depth-one effective twists land exactly there.
\end{proposition}

\begin{proposition}[KK.U and KK.P]\label{prop:kku}
\textup{(Universality)} For any admissible $C^2$ unit $f$ in place of the Gaussian, the RR
asymptotic holds with constant $\log\rho(m) + \log(2\!\int_0^\infty\! f)$: the place
structure (the cone, the wall, the per-place product) is independent of the unit; the
Gaussian's only role in the density theorem is the normalization $2\!\int f = 1$.
\textup{(Fine structure)} At one twisted place, the renormalized defect is $e^{-a}$ times an
explicit $\log p^d$-periodic comb --- not a monotone decay; the comb was masked in earlier
grid measurements by the $e^{-a}$ envelope and is recovered at coefficient level.
\end{proposition}

Both are consumed later: universality by the terminal specification (the Gaussian must be a
\emph{fixed point}, not an import, \S\ref{sec:terminal}); the comb by the breakage analysis
of \S\ref{subsec:nn}.

\subsection{The square dual exists as a duality-renormalized count; the inequality is posed}
\label{subsec:ll}

\begin{theorem}[LL.1; \grade{Theorem-candidate as recorded: classical proof (telescoping,
per-lattice Poisson, doubly-exponential block bounds) with every clause witnessed exactly;
witness \texttt{prime\_\allowbreak int\_\allowbreak ll\_\allowbreak square\_\allowbreak dual}}]\label{thm:ll}
Let $T$ be a finite set of primes, $D_T(a)$ the strict-square effective twist ($m_p = -1$,
$p \in T$), $\rho_T = \prod_p p/(p+1)$. Then: \textup{(i)} the K2-paired resummation of the
divergent dual converges absolutely --- each block's continuum zero modes cancel exactly, and
the unpaired and point-grouped readings do not converge absolutely (\emph{pairing $=$
existence}, the GG mechanism verbatim); \textup{(ii)} the dual object is the explicit signed
measure $c_T(\xi) = \prod_p c_p(v_p(\xi))$; \textup{(iii)} exact inhomogeneous RR:
$\Theta_T(t) = 1 + \tfrac{1}{t}\Theta_T^{\mathrm{dual}}(1/t)$, i.e.\
$h^0(D) - h^0_{\mathrm{ren}}(K_{\mathrm{sq}} - D) = \deg D + \varepsilon(D)$ with
$\varepsilon$ explicit, $0 < \varepsilon = O(e^{-\deg D})$; \textup{(iv)}
$\Theta_T^{\mathrm{dual}}(u) \to \rho_T$ from below with $\log p^2$-periodic fine structure.
\end{theorem}

The index inequality is thereby \emph{posable} for the first time with every term defined:
\[
h^0(D) + h^0_{\mathrm{ren}}(K_{\mathrm{sq}} - D) \;\ge\; \chi_{\mathrm{sq}}(D),
\qquad \chi_{\mathrm{sq}}(D) = \deg D + \varepsilon(D),
\]
identically equivalent to $h^0_{\mathrm{ren}}(K_{\mathrm{sq}} - D) \ge 0$. Its status, as
measured: \emph{false} for the bare renormalized count (sign-indefinite at fine-structure
order on the effective sector; doubly-exponentially false on the Schur-negative sector), and
the entire failure is \emph{one atom}: the dual side of the functional equation lacks the
unit class, whose Fourier dual is the continuum unit. Inserted by hand, the unit repairs the
inequality on the strict family uniformly --- \emph{with zero slack at deep effectivity}: the
unit must couple, not pad.

\subsection{The graph-class test: wrong face, and the breakage map is the comb}
\label{subsec:nn}

The cheap fork before any multi-year commitment: test the hand-inserted unit on the
Frobenius-graph classes, where the zeros' EF weight actually lives. Verdict, deterministic
across runs, breakage arb-certified, freeze protocol verbatim
\grade{Certified/Census; witness \texttt{prime\_\allowbreak nn\_\allowbreak graph\_\allowbreak class\_\allowbreak index}}:
\begin{itemize}
\item the $\Lam(p)$ prime-atom classes ($k = 1$ steps) \emph{fail always}:
$m_+ = -e^{-a}/2\,(1 + o(1))$, certified negative at all tested scales, place-independent,
smooth in $a$ --- the FE-stripped dual carries an exact continuum hole $-1/u$, and an
inserted \emph{atom} cannot cancel a \emph{continuum} deficit;
\item the $(p; 2)$ ladder classes fail recurrently for every $p \ge 5$, onset exactly at
window entry $a = 2\log p$, with negative-phase fraction per scale period
\begin{center}
\small
\begin{tabular}{lcccccccc}
\toprule
$p$ & $2$ & $3$ & $5$ & $7$ & $11$ & $13$ & $17$ & ($d{=}1$ limit) \\
\midrule
neg.\ fraction & $0.00$ & $0.00$ & $0.25$ & $0.30$ & $0.35$ & $0.36$ & $0.38$ & $1.00$ \\
\bottomrule
\end{tabular}
\end{center}
--- the strict-family success of \S\ref{subsec:ll} was a \emph{small-place artifact}
($p \in \{2, 3\}$ are the safe cells; the first breaking cell is $(5; 2)$); the sampled
interior of the \S\ref{subsec:hh} cone fails $5/5$;
\item the margin-vs-$\tau$ identification \emph{fails} (Spearman $-0.499$): the
strict-family inequality is not the dual-cone slack in disguise;
\item the breakage is $\log p^2$-periodic and comb-determined; against the frozen zero data,
$z < 3$ everywhere: \emph{no zero correlation}.
\end{itemize}
Positive yield en route: the functional equation of Theorem \ref{thm:ll} extends to the whole
integer graph dictionary (mixed depths, $4.9\cdot10^{-88}$). The gap specification sharpens
to its final form: the coupled unit is not one object inserted once; it must be
\emph{class- and scale-coupled} --- one continuum unit per place, paid at each window-entry
phase, largest exactly on the classes carrying $\Lam(p)$ --- which is, of course, exactly
what the explicit formula says the archimedean term is.

\subsection{The face hunt: one survivor}\label{subsec:hunt}

Since the lattice dictionary is not unique, the kill was tested for face-dependence. The
scoreboard \grade{Certified per row; witnesses \texttt{prime\_\allowbreak rr\_\allowbreak face\_\allowbreak cycle},
\texttt{prime\_\allowbreak qq\_\allowbreak native\_\allowbreak face\_\allowbreak index}, \texttt{prime\_\allowbreak pp\_\allowbreak ku\_\allowbreak face\_\allowbreak index}}:

\begin{center}
\small
\begin{tabular}{>{\raggedright\arraybackslash}p{0.20\textwidth} p{0.12\textwidth} >{\raggedright\arraybackslash}p{0.56\textwidth}}
\toprule
\textbf{face} & \textbf{verdict} & \textbf{decisive fact} \\
\midrule
F1: GG re-embedding & dead & repairs $d = 1$ exactly; comb-invisible on the census ---
zero signs moved \\[2pt]
F2: theta aggregate & dead & per-place linear program exactly infeasible; the breaking slot
is an exposed PSD-Toeplitz boundary direction \\[2pt]
F3: $h^1$ loophole & dead & deficits strictly inside $(0, \log 2)$, place-independent,
non-additive --- not counts \\[2pt]
F5: native pairing & dead & the hole is the mirror branch's own pole at $s = 0$; nothing was
ever stripped; yield: the gap object $\Delta_p$ with period-mean $\times \log p$ equal to
$\gamma + \log\pi$ ($1.5\cdot10^{-36}$) \\[2pt]
F6--F8 & parked/dead & by citation at enumeration stage \\[2pt]
F4: the $\ku$ cell comb & \textbf{survives} & excludes residues $\{0, -1\}$ mod $q$; the
excluded $0$-coset's Poisson zero-mode is \emph{one continuum unit} --- the FE-homogenizer
supplied internally; killer cells sign-flip as hand-predicted ($+6.47\cdot10^{-3}$ where the
graph-class test certified $-2.21\cdot10^{-3}$); 527-point census, zero negatives \\
\bottomrule
\end{tabular}
\end{center} The
cross-cutting quantification: the repaired \emph{global} aggregate holds while every
\emph{per-place} measure is exactly infeasible --- the coupled unit \emph{threads} places,
and the $\ku$ face is exactly the face whose cells do that threading internally (the
fracture square's internalized sum over places, \S\ref{subsec:ku}). The hunt's negatives and
its survivor explain each other.

\subsection{The survivor becomes theorems: SS.1--SS.5}\label{subsec:ss}

The stress round turned the surviving census into a closed theorem package, with no named
gap and no zeros consumed anywhere \grade{Theorem; witness
\texttt{prime\_\allowbreak ss\_\allowbreak ku\_\allowbreak inequality\_\allowbreak proof}, all proof-internal constants gated}. Write
$S_q = \{w \in \Z : w \bmod q \notin \{0, q-1\}\}$,
$\Theta_q(t) = \sum_{w \in S_q} e^{-\pi w^2 t^2}$, $\rho_q = (q-2)/q$ (the comb
density),
\[
D_q(u) \;=\; \theta(u) - \frac1q \sum_{k \in \Z}
\Bigl(1 + \cos\frac{2\pi k}{q}\Bigr) e^{-\pi k^2 u^2/q^2},
\qquad \theta(u) = \sum_{w \in \Z} e^{-\pi w^2 u^2},
\]
$T(x) = \sum_{m \ge 1} e^{-\pi m^2 x^2}$, $G(x) = 4xT(x)$, and let
\[
m_+(q, a) \;=\; \log\bigl(1 + e^a D\bigr) + \log(1 + D) - a -
\log\bigl(\rho_q (1 + \rho_q)\bigr), \qquad D = D_q(e^a),
\]
be the cell margin (at the calibration inherited unchanged from the graph-class test, for
comparability), with admissibility $a \ge \log(q+1)$ ($u = e^a \ge q+1$).

\begin{theorem}[SS.1, the homogeneous functional equation]\label{thm:ss1}
For every integer $q \ge 2$ and $t > 0$: $\Theta_q(t) = \tfrac{1}{t} D_q(1/t)$ with
\emph{no} inhomogeneous term; every correction is Gaussian; the hole coefficient is
$c_{\mathrm{hole}}(q) = 0$ exactly. (Contrast the half-line comb, whose Poisson carries the
boundary term $-\tfrac12$: the $d = 1$ hole of \S\ref{subsec:nn} carries exactly minus one
continuum unit.) The cell comb equals the two-coset comb iff $q$ is a prime power
(Lemma A, an iff with a CRT counterexample at every composite non-prime-power modulus ---
the unique point where the prime-power hypothesis is consumed).
\end{theorem}

\begin{proof}[Proof sketch]
For the comb identification (Lemma A): since $\gcd(w, w+1) = 1$, a prime $p$ divides at most
one of $w, w+1$, so for $q = p^j$, $q \mid w(w+1)$ iff $q \mid w$ or $q \mid w+1$ --- the
cell comb is the two-coset comb. Conversely, for $q = q_1 q_2$ with coprime
$q_1, q_2 > 1$, CRT produces $w \equiv 0 \ (q_1)$, $w \equiv -1\ (q_2)$ with
$q \mid w(w+1)$ but $w \not\equiv 0, -1 \ (q)$. For the functional equation: each excluded
coset is a \emph{full} lattice coset, so per-coset Poisson summation (proved in the record
from scratch by differentiation under the integral and periodization --- no contour, no
inherited FE) gives $\Theta_q(t) = \tfrac1t D_q(1/t)$ with no leftover term of any other
homogeneity. The contrast is exact: the half-line comb $\{n \ge 1\}$ is \emph{not} a full
coset, and its Poisson carries the boundary term $-\tfrac12$ --- the stripped unit of the
graph-class kill. The zero-mode ledger inside the proof: an \emph{atom} subtraction dualizes
to a continuum hole; a full \emph{coset} subtraction dualizes to a density reduction. That
is the mechanism of the face's survival, made exact.
\end{proof}

\begin{theorem}[SS.2--SS.3, the all-$(q,a)$ inequality and the deep law]\label{thm:ss23}
For every prime power $q \ge 3$ and every admissible $a$:
$m_+(q, a) \ge \log\bigl(1 + e^{-a}/(\rho_q(1+\rho_q))\bigr) \ge \tfrac49 e^{-a} > 0$, with
entry floor $m_+(q, \log(q+1)) \ge 1/(3q)$ and
$q\, m_+(q, \mathrm{entry}) \to 1 - 6T(1) = 0.740696\ldots$; and
$\bigl|e^a \rho_q\, m_+(q, a) - 1\bigr| \le 1.72\,G(e^a/q) + 0.842\, e^{-a}/\rho_q$ ---
one monotone Gaussian-tail bound covers entry to deep; there is no hard middle range.
\end{theorem}

The measured census sits exactly on the proved laws \grade{Certified; same witness}: the
entry products $q\,m_+(q, \mathrm{entry})$ descend monotonically to the derived limit, and
the deep normalized margins approach the zero-slack value $1$ from below, ordered in
$\rho_q$ exactly as the second-order term predicts:
\begin{center}
\small
\begin{tabular}{lccccc}
\toprule
$q$ & $3$ & $9$ & $25$ & $121$ & $169$ \\
\midrule
$q\,m_+(q,\mathrm{entry})$ & $1.6696$ & $0.9638$ & $0.8175$ & $0.7561$ & $0.7517$ \\
\bottomrule
\end{tabular}
\qquad
\begin{tabular}{lccccc}
\toprule
$q$ & $3$ & $4$ & $5$ & $9$ & $25$ \\
\midrule
$e^{8}\rho_q\, m_+(q, 8)$ & $0.99950$ & $0.99966$ & $0.99972$ & $0.99978$ & $0.99982$ \\
\bottomrule
\end{tabular}
\end{center}
(limit $1 - \tfrac32 G(1) = 0.740696\ldots$ on the left; the zero-slack limit $1$ on the
right).

\begin{theorem}[SS.4, evasion not surplus]\label{thm:ss4}
$0 < m_+(q, a) \le 1.298\, e^{-a}/\rho_q$ for all admissible $(q, a)$; the infimum over the
admissible family is $0$; and if $m_+(q, a) \ge (1+\beta)e^{-a}/\rho_q$ persistently for some
$\beta > 0$, then the hole coefficient satisfies $c_{\mathrm{hole}} \ge \beta' > 0$ ---
forbidden by Theorem \ref{thm:ss1} as an identity. \emph{A surplus is equivalent to a
positive hole coefficient. The face evades the kill; it cannot manufacture slack.} The only
surplus route is aggregation across cells, whose assembled margin is $\Theta(1)$ iff
$\sum_{q \le x} \log p(q) \sim x$ --- Chebyshev/PNT-grade, external, flagged with no decision
weight.
\end{theorem}

(SS.5 is the exact repayment-law identity with $O(q^{-2})$ relative error, derived from the
same master identity.) The truncated (``honest'', lattice-completion-free) face was then
settled by the \emph{zero-charge ledger} \grade{Certified;
\texttt{prime\_\allowbreak tt\_\allowbreak truncated\_\allowbreak face\_\allowbreak chi}}: the entry band's continuum charge is
$M_0(q) = (2q - 4) - 2q\rho_q = 0$ exactly (Fractions) --- zero continuum units, against the
$d = 1$ hole's exactly $-1$ --- so the band is chi bookkeeping, absorbable by a derived
zero-parameter correction $\chi^{\mathrm{corr}} = a + \log(\rho_q E(1 + \rho_q E))$,
$E = \erfc(\sqrt{\pi}\, q e^{-a})$; on the corrected chi the honest face passes the full
527-point census with all signs arb-resolved and zero negatives. Bookkeeping can absorb a
zero-charge defect and cannot manufacture a unit: the survivor and the kills are two sides of
one ledger.

\subsection{The depth wall and the second evaluation}\label{subsec:uu}

The remaining stress item was the actual index content: the two-dimensional
weight-$\times$-scale Castelnuovo posing on the $\ku$ face. Verdict
\grade{Certified, full-value, fail-closed; witness \texttt{prime\_\allowbreak uu\_\allowbreak ku\_\allowbreak castelnuovo\_\allowbreak 2d}}:
the 2D form exists, is well-posed, and is Lorentzian everywhere admissible --- and it is
\emph{supply-grade}: it carries no index-grade spectral content, for an exactly characterized
reason. Every entry of every posing is an explicit function of
$(\Lam_p, \log q, i, m_+(q,a))$, all congruence/entry-determined, with no zero contact
($z < 3$ throughout, frozen comparisons). The quantitative mechanism realizes the registered
risk: the deep law of Theorem \ref{thm:ss23} makes the weight diagonal degenerate onto the
light cone, and the discriminant inequality then demands Frobenius depth
$\rho_q \Lam_p^2\, e^a$ against a window that supplies $a/\log q$: \emph{the quadratic step
fails by $e^a$ versus $a$} --- an exponential wall, not a constant. No re-posing that keeps
the margin as weight budget and window admissibility of the scale family can cross it.

The characterization that closes the attack: the entry $B' = \Lam_p - m_+$ occupies
precisely Weil's fixed-point slot ($\Delta \cdot \Gamma_{F^i} = \#\mathrm{Fix}(F^i)$) ---
but Weil extracts RH because $\#\mathrm{Fix}$ has a \emph{second, independent, spectral
evaluation} (Lefschetz), and every $\ku$-face entry has exactly one, the supply one.
Spectral grade would require an independent trace/fixed-point evaluation of the
weight-class-against-Frobenius pairing --- precisely the missing timelike unit at infinity.
The 2D posing does not create that object; it relocates its absence into the cross term of a
quadratic form. No face evades the unit; the face hunt cannot produce the second evaluation;
only a construction can.

\section{The terminal statement}\label{sec:terminal}

\subsection{One object, demanded six times}

Every route in Sections \ref{sec:compact}--\ref{sec:index} terminates at the same missing
object, each time from an independent direction and with an independent clause of its
specification:

\begin{center}
\small
\begin{tabular}{p{0.12\textwidth} p{0.80\textwidth}}
\toprule
\textbf{source} & \textbf{demand} \\
\midrule
BB (\ref{prop:bb1}) & the comb $\leftrightarrow$ continuum coupling that reads the scale
$\kappa = \pi$ (the archimedean Poisson exchange) \\
GG & one continuum unit per place, subtracted, as the existence condition of the global
trace \\
FF (\ref{prop:ff1}) & a rational/connected (Anderson-layer) carrier, which torsion provably
lacks \\
HH (\ref{thm:hh1}) & the signed archimedean coordinate that \emph{jointly cancels} --- the
timelike unit the certified lattice geometry lacks \\
LL (\ref{thm:ll}) & the element making the square dual's functional equation homogeneous:
the unit class whose Fourier dual is the continuum unit \\
UU (\ref{subsec:uu}) & a second, independent evaluation of the weight-against-Frobenius
pairing --- Weil's Lefschetz slot \\
\bottomrule
\end{tabular}
\end{center}

The graph-class test (\S\ref{subsec:nn}) sharpened the binding: the unit must be class- and
scale-coupled --- one continuum unit per place, paid at each window-entry phase, at the
$e^{-a}/2$ scale, in phase with the combs, largest on the classes carrying $\Lam(p)$. The
zero-slack finding of \S\ref{subsec:ll} and the evasion theorem \ref{thm:ss4} together say it
cannot be inserted, only derived: any padding either fails at finite scale (atom against
continuum) or is forbidden as an identity (positive hole coefficient).

\subsection{The computed shadow}

The object has a computed numerical shadow \grade{Certified;
\texttt{prime\_\allowbreak qq\_\allowbreak native\_\allowbreak face\_\allowbreak index}}. On the native face, the single-place gap
$\Delta_p(a) := \sum_{m \in \Z} d(e^{-a}p^m)$ (the Gaussian-versus-partition deficit)
converges, is exactly $\log p$-periodic in $a$ with maxima at window-entry phases, and its
period-mean is $C^*/\log p$ with
\[
C^* \;=\; \int_\R d(e^u)\, du \;=\; \gamma + \log\pi \;=\; 1.72200\ldots
\]
(by the classical Frullani/gamma integral; certified to $1.5\cdot10^{-36}$): the native
smoothed gap per place is \emph{the explicit formula's archimedean constant per $\log p$ of
scale}, emerging from the deficit with no explicit-formula input. The renormalized
$\chi_{\mathrm{ren}}(C; a) = |C| - \Delta_p(a)$ repairs all single-place classes including
the killers by termwise subset monotonicity --- recorded as a \grade{Reading}: it is the
measured shape of the gap, not a forced identity.

\subsection{The specification}

The construction home, scoped against the literature, is a ``liquid'' (condensed/analytic
\cite{ClausenScholze}) category $\mathsf{L}$ with six demanded clauses, abbreviated:
\textbf{LS1} the program's banked carriers embed conservatively (the integral THH tower with
its Anderson fiber sequence, the mod-$p$ towers including $\Xbar$, the strict count layer);
\textbf{LS2} internal $\R$ and $T = \R/\Z$ with $\Aut_{\mathsf L}(T) = \{\pm 1\}$
($2\pi$-rigidity carried inside $\mathsf{L}$), completing the arithmetic fracture square's
missing corner; \textbf{LS3} an internal Fourier self-duality with $F \circ F = (-1)^*$,
normalization torsor pinned to $\{\pm 1\}$ by LS2, and \emph{the Gaussian as fixed unit} ---
$e^{-\pi x^2}$, and with it $\kappa = \pi$, coming out as fixed-point datum, never inserted;
\textbf{LS4} a trace functional inducing, per place, the comb $\leftrightarrow$ continuum
exchange with GG's one-unit-per-place subtraction as its regularization; \textbf{LS5}
compatibility: the trace recovers the van der Geer--Schoof count on the curve shadow and a
finite renormalized dual exactly where the strict family certifiably diverges (making
$K_{\mathrm{sq}}$ structural); \textbf{LS6} zero free gluing choices, the Epstein/DH carrier
guards at the liquid level, and the F-a guard (no zeta continuation, no zero-side input,
explicit formula as comparison only). The scoping facts \grade{from the program record}:
liquid $\R$ and liquid Serre duality/HRR are in print; there are \emph{zero}
Fourier/Poisson statements in the surveyed liquid corpus (verified absence); and
dualizable-category traces are scale-blind by theorem (Efimov \cite{Efimov24}) --- Proposition
\ref{prop:ff1} one level up, citable: the coupling cannot come from trace formalism alone.
To LS1--LS6 the depth wall (\S\ref{subsec:uu}) adds the terminal demand: the coupled unit
must supply an \emph{independent second evaluation} of the weight-against-Frobenius pairing,
of Lefschetz type, growing with depth.

\subsection{What a construction would and would not give}

If an object satisfying LS1--LS6 plus the second-evaluation demand exists, then by the chain
of Sections \ref{sec:index}: $h^0_{\mathrm{ren}} \ge 0$ uniformly on the posed family, the
index inequality holds there, and the program's conserved obstruction is discharged
\emph{at the layer where it now sits}. We state the guard with the same prominence as the
target. By Corollary \ref{cor:feblind} (FE-blindness), achieving the symmetry layer --- all
local dualities, the global functional equation, even the full self-duality of $\Xbar$ ---
is structurally insufficient for RH-grade content: that layer is provably blind to the
constant the explicit formula needs, and Epstein-shaped objects pass every test in it. The
RH-grade content enters only if the index/density layer does work: the growing observable
$\log N(D)$ against the Frobenius family, with the second evaluation closing the Lefschetz
loop. And the charter's failure mode remains live: it is an open possibility that the
coupled unit's existence, in the form demanded, is itself equivalent to zero-repulsion ---
in which case the conservation law wins and the construction project inherits RH rather than
resolving it. The zero-slack findings are exactly the signature such an equivalence would
leave. Nothing in this paper decides that question; the specification is now precise enough
that it can be decided.

\appendix

\section{The no-go table}\label{app:nogo}

Each row is a route that would, if valid, simplify or replace the missing object; each is
closed by a theorem, a certificate, or a citation-grade structural fact, with the witness
script (in \texttt{scripts/\allowbreak numerics/\allowbreak witness/}) named where one
exists.

\begin{center}
\footnotesize
\begin{tabular}{>{\raggedright\arraybackslash}p{0.44\textwidth} >{\raggedright\arraybackslash}p{0.50\textwidth}}
\toprule
\textbf{route} & \textbf{kill (grade; witness)} \\
\midrule
Counts $+$ positivity $+$ self-duality $+$ temperedness force the critical line &
Epstein boundary: all properties hold, zeros off-line (classical; \S\ref{sec:grammar}) \\[2pt]
Integer signature on curve-level (Hall/module) data &
Euler form vanishes identically, $\#\Ext^1 = \#\Hom$ (theorem); Hall forms positive definite
(\texttt{count\_\allowbreak hall\_\allowbreak signature}) \\[2pt]
Atoms at non-prime-powers, or deformed weights (DH-type supply) &
no carrier: residue fields of size $n$ exist iff $n$ prime (theorem;
\texttt{prime\_\allowbreak int\_\allowbreak thh\_\allowbreak supply}) \\[2pt]
KMS-$\tfrac12$ as a branch/argument convention &
all windowed schemes have integer vacuum slopes; naturality demands $\tfrac12$ (Theorem
\ref{thm:rigidity}; \texttt{prime\_\allowbreak int\_\allowbreak hr3\_\allowbreak scheme\_\allowbreak rigidity}) \\[2pt]
Conjugation/Schwarz symmetry as scheme selector &
selects nothing: every real-$c$ scheme passes (theorem; same witness) \\[2pt]
TP of the square as compactification &
BMS triple: $p$-complete, $\varphi$-semilinear, empty archimedean fracture corner
(citation \cite{BMS18}) \\[2pt]
$C_2$-Tate as the archimedean fold carrier &
$(H\F_p)^{tC_2} = 0$ at odd $p$; the fold lives on fixed points (citation \cite{QS21,DMPR21};
\texttt{prime\_\allowbreak int\_\allowbreak thr\_\allowbreak fold}) \\[2pt]
THR's geometric $C_2$ supplies negation directly &
the residual Weyl action shears, never negates (certified;
\texttt{prime\_\allowbreak int\_\allowbreak u\_\allowbreak f1\_\allowbreak probes}) \\[2pt]
Bare Adams family $\{\psi^1..\psi^n\}$ carries the index &
not Lorentzian at any $n \ge 4$ (certified;
\texttt{prime\_\allowbreak int\_\allowbreak adams\_\allowbreak family\_\allowbreak signature}) \\[2pt]
Any invertible (group-like) family carries Castelnuovo growth &
all graphs bidegree $(1,1)$; groups carry no degrees (theorem in the shadow;
\S\ref{subsec:maxgram}) \\[2pt]
Local Euler supply inside the twisted Adams family &
cooperation lattice monomial through $w \le p^2 - 2$ (citation \cite{Kane81,BakerRichter05}) \\[2pt]
Strict (non-derived) grading bridge to the B\"okstedt tower &
universal invariant is the monoid ring $\Z[\Z\setminus 0]$, free; the slot is empty
(Theorem V.1 of the record; \texttt{prime\_\allowbreak cc\_\allowbreak grading\_\allowbreak invariant\_\allowbreak kernel}) \\[2pt]
Integer-dimension (tolerance) layer holds $\zeta'(0)$ &
quantization defect has a pole at $s = 0$, residue $-1/(2\log b)$ (theorem;
\texttt{prime\_\allowbreak cc\_\allowbreak quantization\_\allowbreak defect}) \\[2pt]
Plain Brown--Comenetz self-duality of $\Xbar$ &
unsatisfiable, Erd\H{o}s--Kaplansky (Theorem \ref{thm:dd};
\texttt{prime\_\allowbreak int\_\allowbreak dd\_\allowbreak bc\_\allowbreak dual}) \\[2pt]
\bottomrule
\end{tabular}
\end{center}

\begin{center}
\footnotesize
\begin{tabular}{>{\raggedright\arraybackslash}p{0.44\textwidth} >{\raggedright\arraybackslash}p{0.50\textwidth}}
\toprule
\textbf{route} & \textbf{kill (grade; witness)} \\
\midrule
Borel-route identification of the $t$-shadow (N1-b) &
the required bigraded slot vanishes; naturality kills the class (Theorem EE.2;
\texttt{prime\_\allowbreak int\_\allowbreak ee\_\allowbreak dd4\_\allowbreak lin}) \\[2pt]
Torsion-level duality determines the scale &
the continuum is BC-annihilated; $2\pi$-rigidity has no torsion shadow (Prop.\
\ref{prop:ff1}; \texttt{prime\_\allowbreak int\_\allowbreak ff\_\allowbreak kappa\_\allowbreak xbar}) \\[2pt]
FE-level properness produces the archimedean constant &
FE layer exactly $\kappa$-blind (Prop.\ \ref{prop:bb1};
\texttt{prime\_\allowbreak int\_\allowbreak bb\_\allowbreak global\_\allowbreak assembly}) \\[2pt]
Integer signature flow detects zeros &
bounded zero amplitude vs $O(1)$ quantization floor; density tower seen, zeros not
(census $+$ mechanism; \texttt{prime\_\allowbreak int\_\allowbreak ii\_\allowbreak signature\_\allowbreak flow}) \\[2pt]
Hand-inserted unit repairs the index inequality on graph classes &
$\Lam(p)$ classes fail always; $(p;2)$, $p \ge 5$ recurrently; comb-determined breakage
(certified; \texttt{prime\_\allowbreak nn\_\allowbreak graph\_\allowbreak class\_\allowbreak index}) \\[2pt]
Alternative faces: GG re-embed / theta aggregate / $h^1$ loophole / native pairing &
all dead at certificate grade (one repairs $d{=}1$ but moves no sign; one exactly
LP-infeasible; one non-additive; one relocates the hole to its own pole)
(\texttt{prime\_\allowbreak rr\_\allowbreak face\_\allowbreak cycle}, \texttt{prime\_\allowbreak qq\_\allowbreak native\_\allowbreak face\_\allowbreak index}) \\[2pt]
Per-cell surplus on the surviving $\ku$ face &
surplus $\equiv$ positive hole coefficient, forbidden by the homogeneous FE as an identity
(Theorem \ref{thm:ss4}; \texttt{prime\_\allowbreak ss\_\allowbreak ku\_\allowbreak inequality\_\allowbreak proof}) \\[2pt]
Truncation band reopens the hole &
band charge exactly $0$ continuum units: bookkeeping, not a unit (certified;
\texttt{prime\_\allowbreak tt\_\allowbreak truncated\_\allowbreak face\_\allowbreak chi}) \\[2pt]
2D Castelnuovo posing on the $\ku$ face yields spectral content &
depth wall: $\rho_q\Lam_p^2 e^{a}$ demanded vs $a/\log q$ supplied; single-evaluation
characterization (certified; \texttt{prime\_\allowbreak uu\_\allowbreak ku\_\allowbreak castelnuovo\_\allowbreak 2d}) \\[2pt]
DH / Epstein survive any of the above &
DH dies at the carrier (no ring; negative weights violate
$\Lam = -\langle d,d\rangle \ge 0$); Epstein keeps exactly the carrier-blind half and dies at
$\mu$ / $v_p$ / product formula / effectivity (multiple witnesses) \\
\bottomrule
\end{tabular}
\end{center}

\section{Reproducibility}\label{app:repro}

\subsection*{Witnesses}
Every certified or census-graded claim in this paper is backed by a witness script in
\texttt{scripts/numerics/witness/} of the program repository, run as
\begin{center}
\texttt{python3 -m scripts.numerics.witness.<name>}
\end{center}
Each witness is \emph{fail-closed}: its gates are pre-registered predicates (written before
the first run, a discipline recorded per worker in the program documents), it exits nonzero
on any gate failure, and proof-internal constants are gated at the values derived in the
proofs, not fitted to the outputs. Exact claims use integer/\texttt{Fraction}/$\F_2$
arithmetic; analytic claims use FLINT/Arb ball arithmetic (typical working precision
300--600 bits), and an inequality is reported certified only when the ball resolves its
sign. The load-bearing witnesses, by section:
\begin{itemize}\raggedright
\item \S\ref{sec:grammar}: \texttt{count\_\allowbreak fingerprint}, \texttt{count\_\allowbreak hall\_\allowbreak signature};
\item \S\ref{sec:square}: \texttt{prime\_\allowbreak int\_\allowbreak derived\_\allowbreak square},
\texttt{prime\_\allowbreak int\_\allowbreak thh\_\allowbreak supply}, \texttt{prime\_\allowbreak int\_\allowbreak ku\_\allowbreak ladder},
\texttt{prime\_\allowbreak int\_\allowbreak genus\_\allowbreak audit};
\item \S\ref{sec:ladders}: \texttt{prime\_\allowbreak int\_\allowbreak tate\_\allowbreak ladder},
\texttt{prime\_\allowbreak int\_\allowbreak deninger\_\allowbreak convention}, \texttt{prime\_\allowbreak int\_\allowbreak thr\_\allowbreak fold},
\texttt{prime\_\allowbreak int\_\allowbreak hr3\_\allowbreak scheme\_\allowbreak rigidity}, \texttt{prime\_\allowbreak int\_\allowbreak thr\_\allowbreak p2\_\allowbreak anomaly},
\texttt{prime\_\allowbreak int\_\allowbreak adams\_\allowbreak family\_\allowbreak signature}, \texttt{prime\_\allowbreak monoid\_\allowbreak frobenius\_\allowbreak shadow};
\item \S\ref{sec:compact}: \texttt{prime\_\allowbreak int\_\allowbreak u\_\allowbreak f1\_\allowbreak probes},
\texttt{prime\_\allowbreak int\_\allowbreak z\_\allowbreak u2\_\allowbreak glue}, \texttt{prime\_\allowbreak int\_\allowbreak gc\_\allowbreak oddp\_\allowbreak glue},
\texttt{prime\_\allowbreak int\_\allowbreak gc\_\allowbreak k3\_\allowbreak barnes\_\allowbreak rank}, \texttt{prime\_\allowbreak int\_\allowbreak gc\_\allowbreak k4\_\allowbreak integer\_\allowbreak layer},
\texttt{prime\_\allowbreak int\_\allowbreak dd\_\allowbreak bc\_\allowbreak dual}, \texttt{prime\_\allowbreak int\_\allowbreak ee\_\allowbreak dd4\_\allowbreak lin},
\texttt{prime\_\allowbreak int\_\allowbreak ff\_\allowbreak kappa\_\allowbreak xbar};
\item \S\ref{sec:global}: \texttt{prime\_\allowbreak int\_\allowbreak bb\_\allowbreak global\_\allowbreak assembly},
\texttt{prime\_\allowbreak cc\_\allowbreak quantization\_\allowbreak defect}, \texttt{prime\_\allowbreak int\_\allowbreak gg\_\allowbreak global\_\allowbreak cancellation},
\texttt{prime\_\allowbreak int\_\allowbreak ii\_\allowbreak signature\_\allowbreak flow};
\item \S\ref{sec:index}: \texttt{prime\_\allowbreak hh\_\allowbreak weil\_\allowbreak index\_\allowbreak factorization},
\texttt{prime\_\allowbreak int\_\allowbreak kk\_\allowbreak l1\_\allowbreak density\_\allowbreak rr}, \texttt{prime\_\allowbreak int\_\allowbreak ll\_\allowbreak square\_\allowbreak dual},
\texttt{prime\_\allowbreak nn\_\allowbreak graph\_\allowbreak class\_\allowbreak index}, \texttt{prime\_\allowbreak rr\_\allowbreak face\_\allowbreak cycle},
\texttt{prime\_\allowbreak qq\_\allowbreak native\_\allowbreak face\_\allowbreak index}, \texttt{prime\_\allowbreak pp\_\allowbreak ku\_\allowbreak face\_\allowbreak index},
\texttt{prime\_\allowbreak ss\_\allowbreak ku\_\allowbreak inequality\_\allowbreak proof}, \texttt{prime\_\allowbreak tt\_\allowbreak truncated\_\allowbreak face\_\allowbreak chi},
\texttt{prime\_\allowbreak uu\_\allowbreak ku\_\allowbreak castelnuovo\_\allowbreak 2d}.
\end{itemize}
Approximately seventy witnesses were produced by the program in total; every one was rerun
by the lead at the landing of its cycle.

\subsection*{The freeze protocol}
Any computation that compares an output against nontrivial zeros of $\zeta$ (used in this
program only as post-hoc comparison data, never as construction input) runs under the freeze
discipline: all construction inputs and derived loci are computed and hashed
(\textsc{sha-256}) \emph{before} the zero data is loaded; the comparison statistic ($z$-score
against shuffled/jittered controls) is computed after; the hash is reported and must be
byte-identical across reruns. The signature-flow, graph-class, and face witnesses all
report their freeze hashes in their output and in the program record.

\subsection*{The degrees-of-freedom discipline (G-INT6)}
Every dictionary claim ships a ledger counting free modeling choices (weight grading, scale,
regularization scheme, ladder start, place dictionary) against independently matched
constraints; a claim without surplus constraints is recorded as calibration. The standing
fitted-dictionary control is the number-ring family of \S\ref{subsec:ku}: a construction
claiming to produce $\zeta$'s data must, fed $\mathcal{O}_K$, produce $\xi_K$'s. Under this
audit the exact integer identities (B\"okstedt tower orders, the $\Lam$ dictionary, the
$\mathrm{lcm}$ ladder, the Lindenstrauss--Madsen orders) pass with zero spent choices; the
dimensionful readings (the genus number, ``the constant \emph{is} $\zeta'(0)$'') are
recorded as calibration, and are so labeled in this paper.

\subsection*{Provenance}
This paper compresses program documents 17--37 (\texttt{research/proof/}) and the worker
records cited inline (\texttt{research/work/}); the corrections recorded there (the
ladder-versus-scheme refutation, the genus demotion, the H-r3 vacuity repair, the
scheme-applied weakening at $p = 2$, and the DD.4 display fix) are reproduced in the text as
corrections, not silently absorbed.

\begin{thebibliography}{30}

\bibitem{Weil52} A.~Weil, \emph{Sur les ``formules explicites'' de la th\'eorie des nombres
premiers}, Meddelanden Fr\aa n Lunds Univ. Mat. Sem. (d\'edi\'e \`a M.~Riesz) (1952), 252--265.

\bibitem{Bombieri00} E.~Bombieri, \emph{Remarks on Weil's quadratic functional in the theory
of prime numbers, I}, Rend. Mat. Acc. Lincei (9) \textbf{11} (2000), 183--233.

\bibitem{Yoshida92} H.~Yoshida, \emph{On Hermitian forms attached to zeta functions}, in:
Zeta Functions in Geometry, Adv. Stud. Pure Math. \textbf{21}, Math. Soc. Japan (1992), 281--325.

\bibitem{CC20} A.~Connes and C.~Consani, \emph{Weil positivity and trace formula: the
archimedean place}, Selecta Math. (N.S.) \textbf{27} (2021), no.~77; arXiv:2006.13771.

\bibitem{CCRR22} A.~Connes and C.~Consani, \emph{Riemann--Roch for $\overline{\mathrm{Spec}\,\Z}$},
arXiv:2205.01391.

\bibitem{CCRR23} A.~Connes and C.~Consani, \emph{Riemann--Roch for the ring $\Z$},
arXiv:2306.00456.

\bibitem{Bokstedt} M.~B\"okstedt, \emph{The topological Hochschild homology of $\Z$ and of
$\Z/p$}, unpublished preprint, Bielefeld (ca.~1985).

\bibitem{Hesselholt16} L.~Hesselholt, \emph{Topological Hochschild homology and the Hasse--Weil
zeta function}, in: An Alpine Bouquet of Algebraic Topology, Contemp. Math. \textbf{708},
Amer. Math. Soc. (2018), 157--180; arXiv:1602.01980.

\bibitem{Deninger91} C.~Deninger, \emph{On the $\Gamma$-factors attached to motives},
Invent. Math. \textbf{104} (1991), 245--261.

\bibitem{Deninger92} C.~Deninger, \emph{Local $L$-factors of motives and regularized
determinants}, Invent. Math. \textbf{107} (1992), 135--150.

\bibitem{DMPR21} E.~Dotto, K.~Moi, I.~Patchkoria and S.~P.~Reeh, \emph{Real topological
Hochschild homology}, J. Eur. Math. Soc. \textbf{23} (2021), 63--152.

\bibitem{QS21} J.~D.~Quigley and J.~Shah, \emph{On the parametrized Tate construction and two
theories of real $p$-cyclotomic spectra}, arXiv:1909.03920; \emph{On the equivalence of two
theories of real cyclotomic spectra}, arXiv:2112.07462.

\bibitem{Lee22} D.~Lee, \emph{Topological Hochschild homology of the connective image-of-$J$
spectrum and of connective complex $K$-theory}, arXiv:2206.02411.

\bibitem{BakerRichter05} A.~Baker and B.~Richter, \emph{On the cooperation algebra of the
connective Adams summand}, Tbilisi Math. J. (2008); arXiv:math/0506455.

\bibitem{Kane81} R.~Kane, \emph{Operations in connective $K$-theory}, Mem. Amer. Math. Soc.
\textbf{34} (1981), no.~254.

\bibitem{vdGS} G.~van der Geer and R.~Schoof, \emph{Effectivity of Arakelov divisors and the
theta divisor of a number field}, Selecta Math. (N.S.) \textbf{6} (2000), 377--398;
arXiv:math/9802121.

\bibitem{LM00} A.~Lindenstrauss and I.~Madsen, \emph{Topological Hochschild homology of number
rings}, Trans. Amer. Math. Soc. \textbf{352} (2000), 2179--2204.

\bibitem{Stonek18} B.~Stonek, \emph{Higher topological Hochschild homology of periodic complex
$K$-theory}, Topology Appl. \textbf{282} (2020); arXiv:1801.00156.

\bibitem{GM95} J.~P.~C.~Greenlees and J.~P.~May, \emph{Generalized Tate cohomology},
Mem. Amer. Math. Soc. \textbf{113} (1995), no.~543.

\bibitem{DGI06} W.~G.~Dwyer, J.~P.~C.~Greenlees and S.~Iyengar, \emph{Duality in algebra and
topology}, Adv. Math. \textbf{200} (2006), 357--402.

\bibitem{GreenleesMeier17} J.~P.~C.~Greenlees and L.~Meier, \emph{Gorenstein duality for real
spectra}, Algebr. Geom. Topol. \textbf{17} (2017), 3547--3619.

\bibitem{Efimov24} A.~Efimov, \emph{K-theory and localizing invariants of large categories},
arXiv:2405.12169.

\bibitem{ClausenScholze} D.~Clausen and P.~Scholze, \emph{Lectures on analytic geometry},
lecture notes, Bonn (2019--2022).

\bibitem{BMS18} B.~Bhatt, M.~Morrow and P.~Scholze, \emph{Topological Hochschild homology and
integral $p$-adic Hodge theory}, Publ. Math. IH\'ES \textbf{129} (2019), 199--310;
arXiv:1802.03261.

\bibitem{Zelevinsky81} A.~Zelevinsky, \emph{Representations of finite classical groups: a Hopf
algebra approach}, Lecture Notes in Math. \textbf{869}, Springer (1981).

\end{thebibliography}

\end{document}
